% Texte allemand modernisé : document anonyme publié sur le web (extraits),
% reprenant peut-être une édition du XXe siècle (à vérifier).
% Texte original : Weimarer Ausgabe, t. 53 (1920), p. 417-552.
% Traduction : nouvelle, établie d'après WA 53. Référence consultée : Martin Luther,
% Des Juifs et de leurs mensonges (1543), édition critique, trad. Johannes Honigmann,
% introduction et notes de Pierre Savy, Honoré Champion, 2015, ISBN 978-2-7453-2847-2.

\documentclass[a4paper,12pt]{article}

\usepackage[a4paper,margin=8mm,bmargin=16mm]{geometry}
\usepackage[german,french]{babel}
\usepackage{fontspec}
\usepackage{microtype}
\usepackage[autostyle]{csquotes}

\setmainfont{Alegreya}

\usepackage{paracol}
\newcommand\chunks[2]{%
\begin{leftcolumn*}\begin{otherlanguage}{german}%
{#1}%
\end{otherlanguage}\end{leftcolumn*}%
\begin{rightcolumn}\begin{otherlanguage}{french}%
{#2}%
\end{otherlanguage}\end{rightcolumn}%
}
% https://tex.stackexchange.com/a/505562/295527

\begin{document}
\thispagestyle{empty}
\setlength{\columnsep}{8mm}
\setlength{\columnseprule}{0.4pt}

\begin{paracol}{2}
  \chunks{\section*{Von den Juden und ihren Lügen}}
         {\section*{Des Juifs et de leurs mensonges}}

% Introduction : WA 53, p. 417-419
  \chunks{\subsection*{Einleitung}}
         {\subsection*{Introduction}}
% WA 53, p. 417
  \chunks{Aber der Teufel ist der Welt Gott, und wo Gottes Wort nicht ist, hat er gut machen, nicht allein bei den Schwachen, sondern auch bei den Starken.}
         {Mais le diable est le dieu du monde, et là où la parole de Dieu n’est pas, il a beau jeu, non seulement auprès des faibles, mais aussi auprès des forts.}
  \chunks{\textelp{}}
         {\textelp{}}
% WA 53, p. 419
  \chunks{So wollen wir nun, unsern Glauben zu stärken, der Juden etliche grobe Torheit in ihrem Glauben und Auslegung der Schrift handeln, weil sie so giftig unsern Glauben lästern; kommts irgendeinem Juden zur Besserung, daß er sich schäme, ists desto besser.}
         {Nous voulons donc, pour fortifier notre foi, traiter de quelques grossières sottises des Juifs dans leur foi et dans leur interprétation de l’Écriture, puisqu’ils outragent notre foi avec tant de venin ; si cela peut amener quelque Juif à s’amender en ayant honte, tant mieux.}
  \chunks{\textelp{}}
         {\textelp{}}
% Partie I : WA 53, p. 419-449
  \chunks{\subsection*{I. Der falsche Ruhm der Juden}}
         {\subsection*{I. Les faux titres de gloire des Juifs}}
% Partie I, point 1 : WA 53, p. 419-427
  \chunks{\subsubsection*{1. Die Geburt}}
         {\subsubsection*{1. La naissance}}
% WA 53, p. 426-427
  \chunks{Was heißt aber \enquote{in Sünden geboren} anders, denn in Gottes Zorn und Verdammnis geboren, daß wir von Natur oder Geburt wegen nicht können Gottes Volk oder Kinder sein und unsere Geburt, Ruhm und Adel, Ehre und Preis kein anderer ist noch sein kann, denn daß wir (wo nicht mehr da sein wird denn die leibliche Geburt) schlechthin verdammte Sünder, Gottes Feinde und in allen Ungnaden sind?}
         {Or, que veut dire \enquote{né dans le péché}, sinon né dans la colère et la damnation de Dieu, de sorte que, par nature ou de naissance, nous ne pouvons être le peuple ni les enfants de Dieu, et que notre naissance, nos titres de gloire et notre noblesse, notre honneur et notre louange ne sont ni ne peuvent être autre chose que d’être (s’il n’y a rien de plus que la naissance charnelle) purement et simplement des pécheurs damnés, des ennemis de Dieu, en pleine disgrâce ?}
  \chunks{\textelp{}}
         {\textelp{}}
% Partie I, point 2 : WA 53, p. 427-439
  \chunks{\subsubsection*{2. Die Beschneidung}}
         {\subsubsection*{2. La circoncision}}
% WA 53, p. 432
  \chunks{Soll nun Gott ihr Gebet und Lob erhören und annehmen, müssen sie fürwahr zuvor solchen lästerlichen, schändlichen, falschen, verlogenen Ruhm und Hochmut aus ihren Schulen, Munde und Herzen tun.}
         {Si Dieu doit exaucer et agréer leur prière et leur louange, il faut assurément qu’ils ôtent d’abord de leurs synagogues, de leur bouche et de leur cœur cette vantardise et cet orgueil blasphématoires, infâmes, faux et mensongers.}
  \chunks{\textelp{}}
         {\textelp{}}
% WA 53, p. 433
  \chunks{Alles ängstliche Seufzen und Sehnen und Hoffen ihres Herzens geht dahin, daß sie einmal mit uns Heiden umgehen könnten, wie sie zur Zeit Esthers in Persien mit den Heiden umgingen. O wie lieb haben sie das Buch Esther, das so fein zu ihrer blutdürstigen, rachgierigen, mörderischen Begier und Hoffnung stimmt! Kein blutdürstigeres und rachgierigeres Volk hat die Sonne je beschienen, als die sich dünken lassen, sie seien darum Gottes Volk, daß sie sollen und müssen die Heiden morden und würgen. Und es ist auch das vornehmste Stück, das sie von ihrem Messias erwarten, er solle die ganze Welt durch ihr Schwert ermorden und umbringen. Wie sie denn im Anfang an uns Christen in aller Welt wohl erwiesen und noch gerne täten, wenn sie es könnten, habens auch oft versucht und darüber auf die Schnauze weidlich geschlagen worden sind.}
         {Tous les soupirs, désirs et espoirs angoissés de leur cœur tendent à ceci : pouvoir un jour en user avec nous, gentils, comme ils en usèrent avec les gentils en Perse au temps d’Esther. Oh, comme ils l’aiment, le livre d’Esther, qui s’accorde si bien à leur convoitise et à leur espérance sanguinaires, vindicatives et meurtrières ! Jamais le soleil n’a éclairé peuple plus sanguinaire et plus vindicatif que ceux qui se figurent être le peuple de Dieu afin de devoir, et de devoir à toute force, assassiner et égorger les gentils. Et c’est aussi le point principal qu’ils attendent de leur Messie : qu’il assassine et fasse périr le monde entier par leur épée. Comme ils l’ont bien montré au commencement sur nous, chrétiens, dans le monde entier, et le feraient volontiers encore s’ils le pouvaient ; ils l’ont d’ailleurs souvent tenté, et ont été pour cela rudement frappés sur le museau.}
  \chunks{\textelp{}}
         {\textelp{}}
% Partie I, point 3 : WA 53, p. 439-446
  \chunks{\subsubsection*{3. Das Gesetz}}
         {\subsubsection*{3. La Loi}}
% WA 53, p. 444
  \chunks{Gleichwie die zehn Gebote nicht darum gegeben sind, daß man sich davon sollte rühmen und darüber stolziglich alle Welt verachten, als wären sie darum heilig und Gottes Volk, sondern darum sind sie gegeben, daß man sie halten und Gott darin gehorsam sein solle, wie Mose und alle Propheten solches aufs heftigste treiben. Nicht soll sich der rühmen, der sie hat (wie gesagt ist von Teufeln und bösen Leuten), sondern der soll sich rühmen, der sie hält. Wer sie aber hat und nicht hält, soll sich schämen und erschrecken, als der gewißlich dadurch verdammt wird.}
         {De même, les Dix Commandements n’ont pas été donnés pour qu’on s’en glorifie et qu’on méprise pour cela orgueilleusement le monde entier, comme s’ils étaient par là saints et peuple de Dieu, mais ils ont été donnés pour qu’on les garde et qu’on obéisse en cela à Dieu, comme Moïse et tous les prophètes y insistent avec la plus grande force. Ce n’est pas celui qui les a qui doit s’en glorifier (comme on l’a dit des diables et des méchants), c’est celui qui les garde qui doit s’en glorifier. Mais celui qui les a et ne les garde pas doit avoir honte et s’effrayer, puisqu’il est par là assurément damné.}
  \chunks{\textelp{}}
         {\textelp{}}
% Partie I, point 4 : WA 53, p. 446-449
  \chunks{\subsubsection*{4. Land, Stadt und Tempel}}
         {\subsubsection*{4. Le pays, la ville et le Temple}}
% WA 53, p. 447-448
  \chunks{Können sie doch jetzt noch nicht ihren unsinnigen, rasenden Ruhm lassen, daß sie Gottes Volk seien, so sie nun bei 1500 Jahren vertrieben, verstört und zu Grund verworfen sind. Noch hoffen sie wieder dahin zu kommen, um ihres eigenen Verdienstes willen, denn sie haben ja dessen keine Verheißung, der sie sich trösten könnten, ohne was sie aus ihrem irrigen Dünkel in die Schrift schmieren.}
         {Aujourd’hui encore, ils ne peuvent renoncer à leur vantardise insensée et délirante d’être le peuple de Dieu, alors qu’ils sont depuis près de 1500 ans chassés, ruinés et rejetés jusqu’au fond. Et pourtant ils espèrent y retourner par leurs propres mérites, car ils n’en ont aucune promesse dont ils puissent se consoler, hormis ce que leur imagination égarée leur fait fourrer dans l’Écriture.}
  \chunks{\textelp{}}
         {\textelp{}}
% WA 53, p. 448
  \chunks{Gleich aller Dinge wie unsere Papisten, Bischöfe, Mönche und Pfaffen samt ihrem Anhang wollen mit Gewalt Gottes Volk und Kirche sein, und Gott solle sie ansehen darum, daß sie getauft sind, den Namen haben und im Regiment sitzen.}
         {Il en va tout à fait de même de nos papistes, évêques, moines et curetons, avec tous leurs partisans : ils veulent à toute force être le peuple et l’Église de Dieu, et Dieu devrait faire cas d’eux parce qu’ils sont baptisés, qu’ils portent le nom et qu’ils siègent au gouvernement.}
  \chunks{Da stehen sie wie ein Fels, wenn hunderttausend Apostel kämen und sprächen: \enquote{Ihr seid darum nicht die Kirche, daß ihr solch Wesen führet oder viel eigen Werk und Gottesdienst, wo ihrs gleich aufs beste machet, sondern müsset an alle dem verzweifeln und bloß lauter an der Gnade und Barmherzigkeit in Christo hangen etc. Wo nicht, so seid ihr des Teufels Hure oder Buben Schule und nicht die Kirche.} Ermorden, verbrennen, verjagen möchten sie solche Apostel wohl, aber daß sie sollten ihnen glauben und ihr eigen Tun fahren lassen, da ist keine Hoffnung, da wird nichts aus.}
         {Là, ils tiennent comme un roc, quand bien même cent mille apôtres viendraient leur dire : \enquote{Vous n’êtes pas l’Église parce que vous menez ce genre de vie ou que vous faites quantité d’œuvres et de cultes de votre cru, quand même vous les feriez du mieux possible ; mais il vous faut désespérer de tout cela et vous attacher purement et uniquement à la grâce et à la miséricorde dans le Christ, etc. Sinon, vous êtes la putain du diable ou une école de canailles, et non l’Église.} De tels apôtres, ils voudraient bien les assassiner, les brûler, les chasser ; mais qu’ils les croient et renoncent à leurs propres œuvres, n’y comptez pas, il n’en sortira rien.}
  \chunks{Eben so tun die Türken mit ihrem Gottesdienst auch, alle Rottengeister auch, und ist alles voll Juden, Türken, Papisten und Rotten, die allesamt wollen Kirche und Gottes Volk sein nach ihrem Dünkel und Ruhm, unangeachtet des rechten einigen Glaubens und Gehorsams göttlicher Gebote, dadurch doch allein Gottes Kinder werden und bleiben.}
         {Les Turcs font de même avec leur culte, et tous les esprits sectaires aussi ; tout est plein de Juifs, de Turcs, de papistes et de sectaires qui, tous ensemble, veulent être l’Église et le peuple de Dieu selon leur imagination et leur vantardise, sans égard pour la seule vraie foi et l’obéissance aux commandements divins, par lesquelles seules pourtant on devient et l’on reste enfant de Dieu.}
  \chunks{\textelp{}}
         {\textelp{}}
% Partie II : WA 53, p. 449-511
  \chunks{\subsection*{II. Der Messias}}
         {\subsection*{II. Le Messie}}
% WA 53, p. 449
  \chunks{Nun wollen wir von dem Hauptstück reden, da sie Gott um den Messias bitten. Hier sind sie erst die rechten Heiligen und frommen Kinder, hier wollen sie wahrlich nicht Lügner noch Lästerer, sondern gewisse Propheten sein, daß der Messias nicht gekommen sei, sondern noch kommen solle.}
         {Parlons maintenant du point capital, à savoir qu’ils demandent à Dieu le Messie. C’est là surtout qu’ils sont les vrais saints et les enfants pieux ; là, ils ne veulent assurément pas être des menteurs ni des blasphémateurs, mais des prophètes certains que le Messie n’est pas venu et qu’il doit encore venir.}
  \chunks{Wer will sie hier strafen, daß sie irren oder fehlen? Wenn alle Engel und Gott selbst auf dem Berge Sinai oder zu Jerusalem im Tempel öffentlich sprächen, daß der Messias längst gekommen und nunmehr seiner nicht zu harren wäre, so müßte Gott selbst der Teufel und alle Engel eitel Teufel sein. So gewiß sind diese heiligsten, wahrhaftigsten Propheten, daß der Messias nicht gekommen sei, sondern noch kommen soll. Viel weniger werden sie uns hören, wie sie bisher getan und noch immerfort tun, so sie doch durch viel feiner gelehrter Leute, auch aus ihrem Geschlecht, so stattlich überwunden sind, daß auch Stein und Holz, wo sie halbe Vernunft hätten, solches müßten bekennen, und sie doch wissentlich wider erkannte Wahrheit wüten.}
         {Qui pourrait ici les reprendre, leur dire qu’ils se trompent ou qu’ils s’égarent ? Si tous les anges et Dieu lui-même déclaraient publiquement, sur le mont Sinaï ou à Jérusalem dans le Temple, que le Messie est venu depuis longtemps et qu’il n’y a plus à l’attendre, alors Dieu lui-même devrait être le diable, et tous les anges rien que des diables. Tant ces très saints, ces très véridiques prophètes sont certains que le Messie n’est pas venu, mais qu’il doit encore venir. Bien moins encore nous écouteront-ils, comme ils l’ont montré jusqu’ici et le montrent toujours, alors qu’ils ont été réfutés si vigoureusement par beaucoup de gens excellents et savants, y compris issus de leur propre lignée, que même les pierres et le bois, s’ils avaient la moitié d’une raison, devraient le reconnaître ; et pourtant ils s’acharnent sciemment contre la vérité reconnue.}
  \chunks{Aber ihre verdammten Rabbinen, so es wohl anders wissen, ihre arme Jugend und gemeinen Mann so mutwilliglich vergiften und von der Wahrheit abwenden. Denn ich halte, wo solche Schrift möchte bei dem gemeinen Mann und Jugend gelesen werden, sie sollten ihre Rabbinen alle steinigen und ärger hassen, denn sie uns Christen tun. Aber die Bösewichter wehren, daß es nicht für sie komme, wie mans treulich meinet. Und wenn ich meine Papisten nicht hätte erfahren, so wäre mirs unglaublich gewesen, daß auf Erden so böse Leute sein sollten, die öffentlicher, erkannter Wahrheit, das ist Gott selbst ins Angesicht, wissentlich sollten widerstreben.}
         {Mais leurs damnés rabbins, qui savent bien qu’il en va autrement, empoisonnent de propos délibéré leur pauvre jeunesse et l’homme du commun, et les détournent de la vérité. Car je pense que, si de tels écrits pouvaient être lus par l’homme du commun et par la jeunesse, ils lapideraient tous leurs rabbins et les haïraient plus encore qu’ils ne nous haïssent, nous chrétiens. Mais les scélérats empêchent que cela parvienne jusqu’à eux, si loyale qu’en soit l’intention. Et si je n’avais pas fait l’expérience de mes papistes, il m’aurait paru incroyable qu’il pût y avoir sur terre des gens si méchants qu’ils s’opposent sciemment à la vérité manifeste et reconnue, c’est-à-dire à Dieu lui-même, en face.}
  \chunks{Denn ich hätte mich solches verstockten Mutes in keines Menschen Herz, sondern allein in des Teufels Herz versehen. Aber nun verwundere ich mich nicht, weder der Türken noch der Juden Blindheit, Härtigkeit, Bosheit, weil ich solches muß sehen in den allerheiligsten Vätern der Kirchen, Papst, Kardinal, Bischöfen. O du schrecklicher Zorn und unbegreiflich Gericht der hohen göttlichen Majestät, wie kannst du so veracht sein bei Menschenkindern? daß wir nicht des Augenblicks schlecht zu Tod vor dir erzittern! Wie gar ein unerträglicher Blick bist du, auch in der allerheiligsten Herzen und Augen, wie man in Mose und den Propheten siehet; aber diese Steinherzen und eisernen Seelen dich so trötzlich verspotten.}
         {Car je ne me serais attendu à un tel endurcissement dans le cœur d’aucun homme, mais seulement dans le cœur du diable. Mais à présent je ne m’étonne plus ni de l’aveuglement, ni de la dureté, ni de la méchanceté des Turcs ou des Juifs, puisqu’il me faut voir la même chose chez les très saints pères de l’Église, pape, cardinal, évêques. Ô terrible colère, ô jugement incompréhensible de la haute majesté divine, comment peux-tu être à ce point méprisée parmi les enfants des hommes ? Comment ne tremblons-nous pas sur-le-champ devant toi, à en mourir ? Quel spectacle insoutenable tu es, même pour les cœurs et les yeux les plus saints, comme on le voit en Moïse et dans les prophètes ; et ces cœurs de pierre, ces âmes de fer se moquent de toi avec tant d’insolence !}
% WA 53, p. 449-450
  \chunks{Doch, ob wir an den Juden vielleicht umsonst arbeiten (denn ich droben gesagt, ich wolle mit ihnen nicht disputieren), so wollen wir doch, unsern Glauben zu stärken und die schwachen Christen zu warnen vor den Juden, zuvörderst Gott zu ehren, ihre unsinnige Narrheit unter uns selbst handeln, zu beweisen, daß wir recht glauben und sie ganz irre gehen in diesem Artikel von dem Messias.}
         {Toutefois, même si nous travaillons peut-être en vain auprès des Juifs (car j’ai dit plus haut que je ne voulais pas disputer avec eux), nous voulons cependant, pour fortifier notre foi et mettre en garde les chrétiens faibles contre les Juifs, et avant tout pour honorer Dieu, traiter entre nous de leur folie insensée, afin de prouver que nous croyons droitement et qu’ils s’égarent complètement dans cet article du Messie.}
  \chunks{\textelp{}}
         {\textelp{}}
% Partie II, point 1 : WA 53, p. 450-462
  \chunks{\subsubsection*{1. 1. Mose 49: das Zepter Judas}}
         {\subsubsection*{1. Genèse 49 : le sceptre de Juda}}
% WA 53, p. 450
  \chunks{Erstlich nehmen wir für uns den Spruch Gen. 49: \enquote{Es wird das Zepter von Juda nicht wegkommen, bis daß der Silo komme, und dem werden die Völker anhangen etc.} Diesen Spruch des heiligen Patriarchen Jakob, den er an seinem letzten Ende geredet, haben die neuen fremden Juden bisher vielfältiglich zermartert und gekreuzigt, wider ihr eigen Gewissen. Denn sie wissen wohl, daß ihr Deuten und Drehen lauter Mutwill und Büberei ist. Und gemahnet mich ihrer Glossen eben wie eines bösen halsstarrigen Weibes, das ihrem Mann widerbelfert und das letzte Wort haben will, ob sie wohl weiß, daß sie unrecht hat. Also meinen die verblendeten Leute auch, es sei genug, wenn sie wider den Text und rechten Verstand bellen und plaudern können, fragen nichts danach, daß sie mutwilliglich lügen. Und acht wohl, sie möchten lieber leiden, daß der Spruch nie geschrieben wäre, denn daß sie ihren Sinn sollten lassen; er tut ihnen das gebrannte Leid an, und können nicht für ihm über.}
         {Premièrement, prenons la parole de Genèse~49 : \enquote{Le sceptre ne sera pas ôté de Juda jusqu’à ce que vienne le Silo, et c’est à lui que s’attacheront les peuples, etc.} Cette parole du saint patriarche Jacob, prononcée à sa dernière heure, les Juifs nouveaux et étrangers l’ont jusqu’ici torturée et crucifiée de bien des manières, contre leur propre conscience. Car ils savent bien que leur façon d’interpréter et de tordre n’est que malice délibérée et coquinerie. Et leurs gloses me font penser tout à fait à une femme méchante et entêtée, qui criaille contre son mari et veut avoir le dernier mot, bien qu’elle sache qu’elle a tort. De même, ces gens aveuglés pensent qu’il suffit de pouvoir aboyer et bavarder contre le texte et le sens juste, et ils ne se soucient pas de mentir délibérément. Et je pense bien qu’ils aimeraient mieux que cette parole n’eût jamais été écrite plutôt que de devoir renoncer à leur sens ; elle leur inflige le pire des tourments, et ils ne peuvent passer outre.}
  \chunks{\textelp{}}
         {\textelp{}}
% WA 53, p. 456
  \chunks{Erstlich diejenigen, so nicht vom Text fallen, sondern bleiben bei dem Text, sonderlich bei dem chaldäischen Text, welchen kein vernünftiger Jude leugnen tar, diese drehen sich also aus: Gottes Verheißung ist wohl gewiß, aber unser Sünde hindern, daß nicht geschehen kann, was verheißen ist, darum harren wir noch auf seine Verheißung, bis wir gebüßet haben etc.}
         {Premièrement, ceux qui ne s’écartent pas du texte, mais s’en tiennent au texte, en particulier au texte chaldéen, qu’aucun Juif raisonnable n’ose nier, se dérobent ainsi : La promesse de Dieu est bien certaine, mais nos péchés empêchent que puisse s’accomplir ce qui est promis ; c’est pourquoi nous attendons encore sa promesse, jusqu’à ce que nous ayons expié, etc.}
  \chunks{\textelp{}}
         {\textelp{}}
% WA 53, p. 457
  \chunks{Wohl ists wahr, wo Gott etwas verheißt mit Bedingung oder Vorbehalt und spricht: So du das tun wirst, so will ich dies tun, da stehet die Verheißung auf unserm Tun. Als da er zu Salomo sagt: \enquote{Werdet ihr meine Gebote halten, so soll dies Haus mein Heiligtum sein, wo nicht, so will ichs zerstören.} Aber solcher Art ist diese Verheißung von Messias nicht. Denn er spricht nicht also: Werdet ihr dies und das tun, so soll Messias kommen, wo nicht, so soll er nicht kommen. Sondern frei schenkt ers dahin und spricht: Messias soll kommen um die Zeit, wenn das Zepter von Juda weg sein wird. Solche Verheißung stehet bloß auf göttlicher Wahrheit und Gnaden, die unser Tun nicht achtet noch ansiehet. Darum ist diese Ausflucht der Juden faul und dazu ganz lästerlich.}
         {Il est bien vrai que, lorsque Dieu promet quelque chose sous condition ou avec réserve, et qu’il dit : Si tu fais cela, je ferai ceci, la promesse repose sur nos actes. Ainsi lorsqu’il dit à Salomon : \enquote{Si vous gardez mes commandements, cette maison sera mon sanctuaire ; sinon, je la détruirai.} Mais cette promesse du Messie n’est pas de cette sorte. Car il ne dit pas : Si vous faites ceci et cela, le Messie viendra ; sinon, il ne viendra pas. Au contraire, il en fait librement don, et dit : Le Messie viendra au temps où le sceptre aura été ôté de Juda. Une telle promesse repose uniquement sur la vérité et la grâce divines, qui ne font ni cas ni attention à nos actes. C’est pourquoi cette échappatoire des Juifs ne vaut rien, et elle est de plus tout à fait blasphématoire.}
  \chunks{\textelp{}}
         {\textelp{}}
% WA 53, p. 461
  \chunks{Also gehets diesem Spruch Jakob, wie dem Spruch Christi zu unser Zeit gangen ist: \enquote{Das ist mein Leib, der für euch gegeben wird.} Da marterten die Schwärmer auch ein jeglich Wort besonders und vielfältig, setzten auch das Hinterst zuvorderst, ehe sie den Text wollten annehmen in rechtem Verstand, wie wir gesehen haben. Es scheinet aber hie wohl, wie die Christen, als Lyra, Raymundus, Burgensis und ihresgleichen, dennoch fleißig gewesen und gern die Juden bekehret hätten, haben sie gejagt von einem Wort aufs ander, wie man die Füchse jagt. Aber, wenn sie lange gejagt, so sind zuletzt die Juden auf ihrem Mutwillen verharret und wissentlich wollen irren und nicht von ihren Rabbinen weichen, so müssen wir sie auch lassen fahren und ihr giftige Lästerung und Lügen nicht achten.}
         {Il en va donc de cette parole de Jacob comme il en est allé, de notre temps, de la parole du Christ : \enquote{Ceci est mon corps, qui est donné pour vous.} Là aussi, les exaltés ont torturé chaque mot séparément et de bien des manières, et ont mis devant ce qui était derrière, plutôt que de vouloir accepter le texte dans son sens juste, comme nous l’avons vu. On voit bien ici, cependant, que les chrétiens, comme Nicolas de Lyre, Raymond Martin, Paul de Burgos et leurs semblables, ont du moins été zélés et auraient volontiers converti les Juifs : ils les ont pourchassés d’un mot à l’autre, comme on chasse les renards. Mais après qu’ils les eurent longtemps pourchassés, les Juifs ont fini par s’obstiner dans leur malice, en voulant s’égarer sciemment et ne pas s’écarter de leurs rabbins ; il nous faut donc, nous aussi, les laisser aller, et ne pas faire cas de leurs blasphèmes et de leurs mensonges venimeux.}
  \chunks{\textelp{}}
         {\textelp{}}
% Partie II, point 3 : WA 53, p. 476-484
  \chunks{\subsubsection*{3. Haggai 2}}
         {\subsubsection*{3. Aggée 2}}
% WA 53, p. 476
  \chunks{Weiter nehmen wir für uns der letzten Propheten einen. Es stehet geschrieben Hagg. 2: \enquote{So spricht der HERR: Es ist noch um ein kleine Zeit, daß ich Himmel und Erden, das Meer und Trocken bewegen werde, ja alle Heiden will ich bewegen. Da soll denn kommen der Heiden Trost, Hemdath. Und ich will dies Haus voll Herrlichkeit machen, spricht der HERR Zebaoth. Denn ich hab Silber und Gold genug, spricht der HERR Zebaoth. Es soll die Herrlichkeit dieses letzten Hauses größer sein, denn des ersten gewesen ist, spricht der HERR Zebaoth. Und ich will Friede geben an diesem Ort, spricht der HERR Zebaoth.}}
         {Prenons ensuite l’un des derniers prophètes. Il est écrit en Aggée~2 : \enquote{Ainsi parle le SEIGNEUR : Encore un peu de temps, et j’ébranlerai le ciel et la terre, la mer et le sec ; oui, j’ébranlerai tous les gentils. Alors viendra la consolation des gentils, \emph{hemdath}. Et je remplirai cette maison de gloire, dit le SEIGNEUR Sabaoth. Car j’ai assez d’argent et d’or, dit le SEIGNEUR Sabaoth. La gloire de cette dernière maison sera plus grande que ne fut celle de la première, dit le SEIGNEUR Sabaoth. Et je donnerai la paix en ce lieu, dit le SEIGNEUR Sabaoth.}}
  \chunks{\textelp{}}
         {\textelp{}}
% WA 53, p. 477
  \chunks{Aber hie drehen sie sich also aus. Weil sie die kurze kleine Zeit nicht können leugnen, nehmen sie das Wort \enquote{Trost aller Heiden}, ebräisch \emph{Hemdath}, für sich und kreuzigens, wie sie droben dem Spruch Jakob getan haben mit dem Worte \emph{Schevet} und \emph{Schilo}, wollen, es solle nicht Messias bedeuten, sondern es solle bedeuten: Gold und Silber aller Heiden. Weil das Wort Hemdath heißt eigentlich nach der Grammatica Begierd oder Lust, als da die Heiden begehren, Lust und Liebe zu haben. Und soll der Text also stehen: Noch über ein kurzes, so soll kommen die Begierd aller Heiden. Was ist das? Was begehren die Heiden? Gold, Silber und Kleinod. Möchtest vielleicht fragen: Warum die Juden solche Glossen hie machen? Will ich dir sagen: Der Odem stinkt ihnen nach der Heiden Gold und Silber, denn kein Volk unter der Sonnen geiziger, denn sie sind, gewest ist, noch sind und immer fort bleiben, wie man siehet an ihrem verfluchten Wucher. Und sich auch trösten, wenn ihr Messias kömmt, soll er aller Welt Gold und Silber nehmen und unter sie teilen. Darum, wo sie können die Schrift auf ihren unsättigen Geiz ziehen, das tun sie mit allem Frevel. Du kannst denken, daß Gott und seine Propheten sonst nichts zu weissagen gewußt haben, denn wie der grundlose Geiz der verdammten Juden mit der Heiden Gold und Silber zu sättigen sei.}
         {Mais ici, ils se dérobent ainsi. Comme ils ne peuvent nier ce temps court et bref, ils s’en prennent au mot \enquote{consolation de tous les gentils}, en hébreu \emph{hemdath}, et le crucifient, comme ils l’ont fait plus haut pour la parole de Jacob avec les mots \emph{schevet} et \emph{schilo} : ils veulent qu’il ne signifie pas le Messie, mais l’or et l’argent de tous les gentils. Car le mot \emph{hemdath} signifie proprement, selon la grammaire, désir ou envie, c’est-à-dire ce que désirent les gentils, ce qu’ils convoitent et aiment. Le texte devrait donc se lire ainsi : Encore un peu de temps, et viendra le désir de tous les gentils. Qu’est-ce à dire ? Que désirent les gentils ? De l’or, de l’argent et des joyaux. Tu demanderas peut-être pourquoi les Juifs font ici de telles gloses. Je vais te le dire : ils ne respirent que l’or et l’argent des gentils, car aucun peuple sous le soleil n’a été ni n’est plus avare qu’eux, et ils le restent toujours, comme on le voit à leur maudite usure. Et ils se consolent aussi en se disant que, quand leur Messie viendra, il prendra l’or et l’argent du monde entier et le partagera entre eux. C’est pourquoi, partout où ils peuvent tirer l’Écriture vers leur avarice insatiable, ils le font avec la dernière impudence. Figure-toi que Dieu et ses prophètes n’auraient rien su prophétiser d’autre que la manière de rassasier l’avarice sans fond des Juifs damnés avec l’or et l’argent des gentils !}
  \chunks{\textelp{}}
         {\textelp{}}
% WA 53, p. 478
  \chunks{Weiter, wie gehets zu, daß solche hochverständige Meister und weise, heilige Propheten nicht auch das Wort Begierd (\emph{Hemdath}) ausbreiten auf alle ander Begierd der Heiden? Denn die Heiden begehren nicht allein Gold und Silber, sondern auch hübsche Fräulein, und die Frauen hübsche Jünglinge. Und wo unter den Heiden nicht eitel Juden (ich wollt sagen Geizwänste) sind, die ihrem eigen Leibe nichts Guts tun, so begehren sie auch schöne Häuser, Gärten, Vieh, Güter, item gute Tage, Kleider, fressen, saufen, tanzen, spielen und allerlei Freude. Warum deuten sie nun nicht diesen Spruch des Propheten, daß solche Begierd aller Heiden auch sollen gen Jerusalem in kurz kommen, auf daß die Juden allein ihren Bauch füllen und sich in der Welt Wollust weiden? Denn solch Wesen verheißt Mahomet seinen Sarazenen, und ist in dem ein rechter Jude, und die Juden rechte Sarazenen, nach dieser Deutung.}
         {Ensuite, comment se fait-il que des maîtres si pleins d’intelligence, de si sages et saints prophètes n’étendent pas aussi le mot désir (\emph{hemdath}) à tous les autres désirs des gentils ? Car les gentils ne désirent pas seulement l’or et l’argent, mais aussi de jolies demoiselles, et les femmes de beaux jeunes gens. Et là où il n’y a pas, parmi les gentils, que des Juifs (je voulais dire : des panses avares)\footnote{Faux lapsus : Luther feint de se reprendre, comme si \enquote{Juifs} lui avait échappé, pour faire du mot un synonyme d’\enquote{avares}. La figure passe en français, mais non le mot \foreignlanguage{german}{\emph{Geizwänste}}, composé injurieux courant au \textsc{xvi}\textsuperscript{e}~siècle, formé sur \foreignlanguage{german}{\emph{Geiz}} (\enquote{avarice}, mais aussi \enquote{cupidité}) et \foreignlanguage{german}{\emph{Wanst}} (\enquote{panse}, \enquote{bedaine}, mot péjoratif pour le ventre) : le ventre avide d’amasser, qui pourtant se refuse tout. Même procédé plus bas (WA~53, p.~479), où le jeu repose cette fois sur la quasi-homophonie et ne peut se traduire.}, qui refusent tout bien-être à leur propre corps, ils désirent aussi de belles maisons, des jardins, du bétail, des biens, et encore du bon temps, des habits, bâfrer, se soûler, danser, jouer et toutes sortes de plaisirs. Pourquoi donc n’interprètent-ils pas cette parole du prophète en ce sens que tous ces désirs des gentils doivent aussi venir sous peu à Jérusalem, afin que les Juifs seuls se remplissent le ventre et se repaissent des voluptés du monde ? Car c’est le genre de vie que Mahomet promet à ses Sarrasins ; en cela il est un vrai Juif, et les Juifs, à suivre cette interprétation, de vrais Sarrasins.}
  \chunks{Noch eine Begierd haben die Heiden. Wie haben doch diese weisen, klugen Deutler derselbigen können vergessen? Das wundert mich. Die Heiden sterben und haben viel Krankheit, Armut, allerlei Not und Angst, und ist keiner, der nicht aufs allerhöchst begehrt, daß er möcht nicht sterben, nicht Not, Elend, Krankheit leiden, oder derselben ja bald los werden und sicher sein. Solche Begierd ist so stark in ihnen, daß sie die andern Begierd alle gern drum geben, wie man täglich siehet in der Erfahrung. Warum deuten sie denn nicht, daß solch Begierd aller Heiden auch kürzlich gen Jerusalem in den Tempel kommen solle? Pfui euch hie, pfui euch dort, und wo ihr seid, ihr verdammten Juden, daß ihr diese ernste, herrliche, tröstliche Wort Gottes so schändlich auf euren sterblichen, madigen Geizwanst ziehen türet, und schämet euch nicht, euren Geiz so gröblich an den Tag zu geben. Seid ihr doch nicht wert, daß ihr die Biblia von außen sollet ansehen, schweige, daß ihr drinnen lesen sollet. Ihr solltet allein die Biblia lesen, die der Sau unter dem Schwanz stehet, und die Buchstaben, so daselbst herausfallen, fressen und saufen; das wäre eine Bibel für solche Propheten, die der göttlichen Majestät Wort, so man mit allen Ehren, Zittern und Freuden hören sollt, so säuisch zu wühlen und so schweinisch zu reißen.}
         {Les gentils ont encore un autre désir. Comment ces sages et habiles faiseurs d’interprétations ont-ils bien pu l’oublier ? Cela m’étonne. Les gentils meurent, et ils ont beaucoup de maladies, de pauvreté, toutes sortes de détresses et d’angoisses ; et il n’est personne qui ne désire par-dessus tout ne pas mourir, ne pas souffrir la détresse, la misère, la maladie, ou du moins en être bientôt délivré et se trouver en sûreté. Ce désir est si fort en eux qu’ils donneraient volontiers pour lui tous les autres désirs, comme on le voit chaque jour par expérience. Pourquoi donc n’interprètent-ils pas ce passage en ce sens que ce désir de tous les gentils doit lui aussi venir sous peu à Jérusalem, dans le Temple ? Honte à vous ici, honte à vous là, et partout où vous êtes, Juifs damnés, qui osez tirer si honteusement ces paroles de Dieu, graves, glorieuses et consolantes, vers votre panse avare\footnote{Même expression (\foreignlanguage{german}{\emph{Geizwanst}}) qu’au paragraphe précédent.}, mortelle et promise aux vers, et qui n’avez pas honte d’étaler si grossièrement votre avarice au grand jour ! Vous n’êtes pas dignes de regarder la Bible du dehors, encore moins d’y lire. Vous ne devriez lire que la Bible qui se trouve sous la queue de la truie\footnote{Allusion à la \foreignlanguage{german}{\emph{Judensau}} (\enquote{truie des Juifs}), bas-relief injurieux de l’église paroissiale de Wittenberg, où Luther prêchait, encore en place aujourd’hui ; Luther le décrit dans \foreignlanguage{german}{\emph{Vom Schem Hamphoras}} (1543).}, et bâfrer et boire les lettres qui en tombent : voilà une Bible pour de tels prophètes, qui fouillent comme des truies et déchirent comme des porcs la parole de la divine Majesté, qu’on devrait écouter avec tout l’honneur, avec tremblement et avec joie.}
  \chunks{\textelp{}}
         {\textelp{}}
% WA 53, p. 479
  \chunks{Weiter, da nun der Prophet sagt: \enquote{Die Herrlichkeit dieses letzten Hauses soll größer sein, denn des ersten}, da laßt uns hören die edlen und beschissen (beschnitten sage ich) Heiligen und weisen Propheten, so uns Christen wollen zu Juden machen. Die große Herrlichkeit des letzten Tempels über den ersten ist diese, daß der letzte, das ist dieser Tempel Haggais, hat zehn Jahr länger gestanden denn der Tempel Salomos etc. Ah, daß sie doch einen guten Astronomen hätten hie gehabt, der es ein wenig genauer hätte abgerechnet! Der würde vielleicht gefunden haben, daß es zu den zehn Jahren noch drei Monden, zwo Wochen, fünf Tage, sieben Stund, zwölf Minuten und zehn halbe Minuten drüber gewesen wären. Wenn etwo ein Kram wäre, der Schamrot feil hätte, dürft ich den Juden ein paar Gulden schenken, daß sie hin gingen und ein Pfund kauften und doch die Stirn, Augen und Backen damit bestrichen, wollten sie ja ihr unverschämt Herz und Zungen nicht damit bestreichen. Oder meinen sie, daß sie mit Stöcken und Klötzen reden, wie sie selbst sind, die ungelehrten groben Esel?}
         {Ensuite, puisque le prophète dit : \enquote{La gloire de cette dernière maison sera plus grande que celle de la première}, écoutons donc les nobles et merdeux (circoncis, veux-je dire)\footnote{Jeu sur \foreignlanguage{german}{\emph{beschissen}} (\enquote{merdeux}, \enquote{souillé}) et \foreignlanguage{german}{\emph{beschnitten}} (\enquote{circoncis}), que la quasi-homophonie rend intraduisible ; même faux lapsus que plus haut (p.~478).} saints et sages prophètes qui veulent faire de nous, chrétiens, des Juifs. La grande gloire du dernier Temple sur le premier, la voici : le dernier, c’est-à-dire ce Temple d’Aggée, est resté debout dix ans de plus que le Temple de Salomon, etc. Ah ! que n’ont-ils eu ici un bon astronome, qui eût fait le compte un peu plus exactement ! Il aurait peut-être trouvé qu’aux dix ans s’ajoutaient encore trois lunes, deux semaines, cinq jours, sept heures, douze minutes et dix demi-minutes. S’il y avait quelque part une boutique où l’on vendît du rouge de la honte, je pourrais bien faire cadeau aux Juifs de quelques florins, pour qu’ils aillent en acheter une livre et s’en enduisent au moins le front, les yeux et les joues, puisqu’ils ne veulent pas en enduire leur cœur et leur langue effrontés. Ou bien pensent-ils parler à des bûches et à des souches, comme eux-mêmes, ces ânes ignares et grossiers ?}
  \chunks{Es waren viel grauer alter Männer und Weiber, auch wohl Bettler und böse Buben zu Jerusalem, da Salomo, ein Jüngling bei zwanzig Jahren, ein herrlicher König ward. Sollten sie darum größer Herrlichkeit haben denn Salomo? Vielleicht ist der Maulesel Davids, darauf Salomo König ward, älter gewesen denn Salomo. Sollt er drum herrlicher sein denn Salomo? Aber also sollen anlaufen, stürzen und fallen, die Gott ohne Unterlaß Lügen strafen und selbs wollen recht haben. Bessers sind sie nicht wert, denn daß sie über die Biblia eitel solche Glossen, Narrenwerk und ihre eigen Schande schreiben (wie sie denn mit großem Fleiß tun). Darum hüte dich, lieber Christ, vor den Juden, die du hieraus siehest, wie sie durch Gottes Zorn dem Teufel übergeben sind, der sie nicht allein des rechten Verstands in der Schrift, sondern auch gemeiner menschlicher Vernunft, Scham und Sinn beraubt hat und treibt durch sie nur eitel Spott mit der heiligen Schrift, daß man hinfort in keinem andern Stück ihnen mehr trauen noch glauben kann, ob ihnen gleich zuweilen ein wahr Wort entfiele. Denn wer so spöttisch und schändlich mit dem schrecklichen Wort Gottes gaukeln tar, wie du hie siehest und droben im Spruch Jakob auch gesehen hast, der muß keinen guten Geist bei sich haben. Darum, wo du einen rechten Juden siehest, magst du mit gutem Gewissen ein Kreuz für dich schlagen und frei sicher sprechen: Da gehet ein leibhaftiger Teufel.}
         {Il y avait à Jérusalem beaucoup de vieillards aux cheveux gris, hommes et femmes, et sans doute aussi des mendiants et des vauriens, quand Salomon, jeune homme d’environ vingt ans, devint un roi glorieux. Devaient-ils pour autant avoir une gloire plus grande que Salomon ? Peut-être le mulet de David, sur lequel Salomon devint roi, était-il plus âgé que Salomon. Devait-il pour cela être plus glorieux que Salomon ? Mais c’est ainsi que doivent se heurter, trébucher et tomber ceux qui, sans relâche, accusent Dieu de mensonge et veulent avoir raison eux-mêmes. Ils ne méritent pas mieux que d’écrire sur la Bible de pures gloses de ce genre, des bouffonneries et leur propre honte (ce qu’ils font d’ailleurs avec grand zèle). Garde-toi donc, cher chrétien, des Juifs : tu vois par là comment la colère de Dieu les a livrés au diable, qui les a privés non seulement du sens juste de l’Écriture, mais aussi de la raison humaine commune, de la pudeur et du bon sens, et qui, par eux, ne fait que se jouer de la sainte Écriture, de sorte qu’on ne peut plus désormais leur faire confiance ni les croire en aucune autre matière, même s’il leur échappait parfois une parole vraie. Car celui qui ose faire le bateleur de façon si frivole et si honteuse avec la terrible parole de Dieu, comme tu le vois ici et comme tu l’as vu plus haut à propos de la parole de Jacob, ne peut avoir en lui aucun bon esprit. C’est pourquoi, partout où tu vois un vrai Juif, tu peux en toute bonne conscience faire le signe de la croix\footnote{Geste de conjuration, comme pour se protéger des mauvais esprits (glose de WA : \foreignlanguage{german}{\emph{wie zum Schutz gegen böse Geister}}).} et dire franchement, à coup sûr : Voilà un diable en chair et en os qui passe.}
  \chunks{\textelp{}}
         {\textelp{}}
% WA 53, p. 480
  \chunks{So ist nun dies der Spruch: \enquote{Noch über ein kleines, so will ich bewegen Himmel und Erden, das Meer und das Trocken (das sind die Inseln im Meer), und soll kommen Hemdath aller Heiden}, das ist Messias, aller Heiden Begierd, welches wir \enquote{Trost} verdeutscht haben. Denn Begierd ist nicht verständlich genug, weil es im Deutschen heißt: die innerlichen Lüst und Begierd im Herzen, \emph{active}. Aber hie heißt es Begierd, das äußerliche, \emph{passive}, des ein Herz begehrt. Wenn mans auch verdeutscht also: Aller Heiden Freude und Wonne, wäre auch nicht falsch. Summa, es ist Messias, an dem die ungläubigen, verstockten Juden eitel Unlust, Ekel und Greuel sollten haben, wie Jesaja 53 von ihnen weissagt. Aber den Heiden sollt er willkommen heißen, als ihres Herzen Freude, Lust, aller Wunsch und Begierd. Denn er bringet ihnen Erlösung von Sünden, Tod, Teufel, Hölle und allem Übel, ewiglich. Ja, das ist der Heiden Begierd, Herzen Lust, Freude und Trost.}
         {Voici donc ce que dit ce passage : \enquote{Encore un peu de temps, et j’ébranlerai le ciel et la terre, la mer et le sec (ce sont les îles dans la mer), et viendra le \emph{hemdath} de tous les gentils}, c’est-à-dire le Messie, le désir de tous les gentils, que nous avons traduit en allemand par \enquote{consolation}\footnote{Luther renvoie à sa traduction de la Bible, où Aggée~2,~7 porte \foreignlanguage{german}{\enquote{der Heiden Trost}}.}. Car \enquote{désir} n’est pas assez clair, puisqu’en allemand le mot désigne les envies et les désirs intérieurs du cœur, au sens actif. Mais ici, désir s’entend de la chose extérieure, au sens passif, que désire un cœur. Si on le traduisait aussi par \enquote{la joie et les délices de tous les gentils}, ce ne serait pas faux non plus. Bref, c’est le Messie, pour qui les Juifs incrédules et endurcis ne devaient avoir que dégoût, nausée et horreur, comme Isaïe~53 le prophétise à leur sujet. Mais pour les gentils, il devait être le bienvenu, comme la joie et le plaisir de leur cœur, l’objet de tous leurs souhaits et de tout leur désir. Car il leur apporte la délivrance du péché, de la mort, du diable, de l’enfer et de tout mal, pour l’éternité. Oui, voilà le désir des gentils, le plaisir de leur cœur, leur joie et leur consolation.}
  \chunks{\textelp{}}
         {\textelp{}}
% WA 53, p. 481
  \chunks{Ja, da steckts, das ist der Hader, da hebt sichs über, das macht die Juden toll und töricht und treibet sie zu solchem verdammten Sinn, daß sie alle Sprüche der Schrift so schändlich verkehren müssen, nämlich: sie wollens nicht, sie könnens nicht leiden, daß wir Heiden sollten ihnen vor Gott gleich sein und der Messias sollt so wohl unser als ihr Trost und Freude sein. Ehe sie das litten, sage ich, daß wir Heiden, die von ihnen ohne Unterlaß verspeiet, vermaledeiet, verflucht, gelästert, geschändet werden, sollten mit ihnen am Messias teil haben, Miterben und ihre Brüder heißen, sie kreuzigten eher noch zehn Messias und schlügen Gott (wo es möglich wäre) selber tot, mit allen Engeln und allen Kreaturen, und sollten sie tausend Höllen für eine verdienen. So gar ein unbegreiflicher harter Stolz steckt in dem edlen Blut der Väter und beschnittenen Heiligen: sie wollen den Messias allein haben und der Welt Herren sein, die verfluchten Gojim sollen Knechte sein, ihr Begierd, das ist ihr Gold und Silber, den Juden geben und sich schlachten lassen wie das arme Vieh.}
         {Oui, c’est là que gît le lièvre, voilà l’objet de la querelle, c’est de là que tout part ; voilà ce qui rend les Juifs fous et insensés, et les pousse à un état d’esprit si damné qu’il leur faut défigurer si honteusement tous les passages de l’Écriture : c’est qu’ils ne veulent pas, qu’ils ne peuvent pas souffrir que nous, gentils, soyons leurs égaux devant Dieu, et que le Messie soit notre consolation et notre joie aussi bien que la leur. Plutôt que de souffrir, dis-je, que nous, gentils, qu’ils ne cessent de conspuer, de maudire, d’exécrer, de blasphémer, de couvrir d’opprobre, ayons part avec eux au Messie, et soyons appelés cohéritiers et leurs frères, ils crucifieraient encore dix Messies et tueraient Dieu lui-même (si c’était possible), avec tous les anges et toutes les créatures, dussent-ils mériter mille enfers au lieu d’un. Tant est incompréhensible et dur l’orgueil qui loge dans le noble sang des pères et des saints circoncis : ils veulent avoir le Messie pour eux seuls et être les maîtres du monde ; les maudits goyim doivent être leurs valets, donner aux Juifs leur désir, c’est-à-dire leur or et leur argent, et se laisser égorger comme le pauvre bétail.}
  \chunks{Ehe sie diesen Sinn lassen, so bleiben sie lieber wissentlich und ewiglich verloren.}
         {Plutôt que de renoncer à cet état d’esprit, ils aiment mieux rester perdus, sciemment et éternellement.}
  \chunks{Sie haben solchen giftigen Haß wider die Gojim von Jugend auf von ihren Eltern und Rabbinern eingesoffen und saufen ihn noch ohne Unterlaß in sich, daß es ihnen, wie der 109. Psalm sagt, durch Blut und Fleisch, durch Mark und Bein gegangen, ganz und gar Natur und Leben geworden ist. Und so wenig sie Fleisch und Blut, Mark und Bein ändern können, so wenig können sie solchen Stolz und Neid ändern, sie müssen so bleiben und verderben, wenn Gott nicht besondere hohe Wunder tut.}
         {Ils se sont gorgés dès leur jeunesse, auprès de leurs parents et de leurs rabbins, d’une haine si venimeuse contre les goyim, et ils continuent de s’en gorger sans relâche, que, comme dit le psaume 109, elle leur est entrée dans le sang et dans la chair, dans la moelle et dans les os, et qu’elle est devenue entièrement leur nature et leur vie. Et aussi peu ils peuvent changer leur chair et leur sang, leur moelle et leurs os, aussi peu peuvent-ils changer un tel orgueil et une telle envie : il leur faut rester ainsi et périr, à moins que Dieu ne fasse des miracles particulièrement grands.}
  \chunks{\textelp{}}
         {\textelp{}}
% WA 53, p. 482
  \chunks{Darum wisse du, lieber Christ, und zweifle nicht daran, daß du nächst dem Teufel keinen bittereren, giftigeren, heftigeren Feind hast als einen rechten Juden, der mit ernst ein Jude sein will. Es mögen wohl unter ihnen sein, die da glauben, was die Kuh oder Gans glaubt, doch hängt ihnen allen das Geblüt und die Beschneidung an. Daher gibt man ihnen oft in den Historien Schuld, daß sie die Brunnen vergiftet, Kinder gestohlen und zerpfriemt haben, wie zu Trient, Weißensee etc. Sie sagen wohl nein dazu. Aber, es sei oder nicht, so weiß ich wohl, daß am vollen, ganzen, bereiten Willen bei ihnen nichts fehlt, wenn sie mit der Tat, heimlich oder offenbar, dazu kommen könnten. Des versieh dich gewißlich und richte dich danach!}
         {Sache donc, cher chrétien, et n’en doute pas : après le diable, tu n’as pas d’ennemi plus amer, plus venimeux, plus violent qu’un vrai Juif, qui veut sérieusement être juif. Il se peut bien qu’il y en ait parmi eux qui croient ce que croit la vache ou l’oie ; mais le sang et la circoncision leur collent à tous à la peau. C’est pourquoi on les accuse souvent, dans les chroniques, d’avoir empoisonné les puits, volé des enfants et percé leur corps à coups d’alène, comme à Trente, à Weissensee, etc. Ils disent bien que non. Mais qu’il en soit ainsi ou non, je sais bien qu’il ne manque rien chez eux à la volonté pleine, entière et prête, s’ils pouvaient en venir à l’acte, en secret ou ouvertement. Attends-toi à cela à coup sûr, et règle-toi là-dessus !}
% WA 53, p. 482
  \chunks{Tun sie aber etwas Gutes, so wisse, daß es nicht aus Liebe noch dir zugute geschieht, sondern weil sie Raum haben müssen, bei uns zu wohnen, müssen sie aus Not etwas tun, aber das Herz bleibt und ist, wie ich gesagt habe.}
         {Mais s’ils font quelque chose de bon, sache que ce n’est ni par amour ni pour ton bien ; c’est parce qu’il leur faut avoir la possibilité d’habiter parmi nous qu’ils sont obligés de faire quelque chose ; mais le cœur reste ce qu’il est, comme je l’ai dit.}
  \chunks{\textelp{}}
         {\textelp{}}
% WA 53, p. 482-483
  \chunks{Es möchte ein Mensch, der den Teufel nicht kennt, sich wohl verwundern, daß sie den Christen vor andern so feind sind, da sie doch nicht Ursache dazu haben, weil wir ihnen alles Gute tun. Sie leben bei uns zu Hause, unter unserm Schutz und Schirm, brauchen Land und Straßen, Markt und Gassen; dazu sitzen die Fürsten und Obrigkeit, schnarchen und haben das Maul offen, lassen die Juden aus ihrem offenen Beutel und Kasten nehmen, stehlen und rauben, was sie wollen, das ist: sie lassen sich selbst und ihre Untertanen durch der Juden Wucher schinden und aussaugen und mit ihrem eigenen Gelde sich zu Bettlern machen. Denn die Juden als Fremdlinge sollten wahrlich und gewißlich nichts haben, und was sie haben, das muß gewißlich unser sein. So arbeiten sie nicht, verdienen uns nichts ab, auch schenken oder geben wir ihnen nichts, dennoch haben sie unser Geld und Gut und sind damit unsere Herren in unserm eigenen Lande und in ihrer Verbannung. Wenn ein Dieb zehn Gulden stiehlt, so muß er hängen, raubt er auf der Straße, so ist der Kopf verloren. Aber wenn ein Jude zehn Tonnen Goldes durch seinen Wucher stiehlt und raubt, so ist er lieber als Gott selbst.}
         {On pourrait bien s’étonner, quand on ne connaît pas le diable, qu’ils soient plus ennemis des chrétiens que de tous les autres, alors qu’ils n’ont aucune raison de l’être, puisque nous leur faisons tout le bien possible. Ils vivent chez nous, sous notre protection et notre sauvegarde, ils usent du pays et des routes, des marchés et des rues ; et par-dessus le marché, les princes et les autorités sont assis là, ronflent et restent bouche bée, et laissent les Juifs puiser dans leur bourse et leur coffre ouverts, voler et piller ce qu’ils veulent ; c’est-à-dire qu’ils se laissent écorcher et sucer, eux et leurs sujets, par l’usure des Juifs, et se réduisent à la mendicité avec leur propre argent. Car les Juifs, en tant qu’étrangers en exil, ne devraient assurément rien avoir, et ce qu’ils ont doit assurément être à nous. Or ils ne travaillent pas, ne gagnent rien à notre service, et nous ne leur faisons pas non plus de cadeaux ni de dons ; pourtant ils ont notre argent et nos biens, et sont par là nos maîtres dans notre propre pays et dans leur exil. Si un voleur vole dix florins, il doit être pendu ; s’il détrousse sur les routes, il y perd la tête. Mais quand un Juif vole et pille dix tonnes d’or par son usure, il est plus cher que Dieu lui-même.}
% WA 53, p. 483
  \chunks{Zum Wahrzeichen: sie rühmen es getrost und stärken ihren Glauben und giftigen Groll wider uns, sprechen untereinander: \enquote{Halte fest, siehe, wie Gott mit uns ist und sein Volk auch in der Verbannung nicht verläßt. Wir arbeiten nicht, haben gute faule Tage, die verfluchten Gojim müssen für uns arbeiten, wir aber kriegen ihr Geld, damit sind wir ihre Herren, sie aber unsere Knechte.}}
         {Et la preuve en est qu’ils s’en vantent hardiment, et qu’ils fortifient leur foi et leur rancune venimeuse contre nous en se disant entre eux : \enquote{Tiens bon, vois comme Dieu est avec nous et n’abandonne pas son peuple, même en exil. Nous ne travaillons pas, nous avons de bonnes journées de paresse ; les maudits goyim doivent travailler pour nous, et nous, nous empochons leur argent ; ainsi nous sommes leurs maîtres, et eux nos valets.}}
  \chunks{\enquote{Halt fest, lieben Kinder Israel, es wird besser werden, unser Messias wird noch kommen, wenn wir so fortfahren und aller Heiden Hemdath durch Wucher und sonst auch an uns gebracht haben.} Wohlan, das leiden wir von ihnen, unter unserm Schutz und Schirm, noch fluchen sie uns, wie gesagt. Aber davon hernach weiter.}
         {\enquote{Tenez bon, chers enfants d’Israël, cela ira mieux : notre Messie viendra encore, quand nous aurons continué ainsi et fait passer entre nos mains, par l’usure et autrement, le \emph{hemdath} de tous les gentils.} Eh bien, voilà ce que nous souffrons d’eux, sous notre protection et notre sauvegarde ; et pourtant ils nous maudissent, comme on l’a dit. Mais nous y reviendrons plus loin.}
  \chunks{\textelp{}}
         {\textelp{}}
% Partie II, point 4 : WA 53, p. 484-511
  \chunks{\subsubsection*{4. Daniel 9: die siebzig Wochen}}
         {\subsubsection*{4. Daniel 9 : les soixante-dix semaines}}
% WA 53, p. 502
  \chunks{Auch wenn sie nicht so stockstarrblind wären, sollte sie ihr eigen, auch das gar grobe äußerliche Leben wohl überzeugen, was sie für Buße tun. Denn sie stecken voller Zauberei, Gaukelei, mit Zeichen, Figuren und des Namens Tetragrammaton (das ist Abgötterei), voll Neides und Stolzes, dazu eitel Diebe und Räuber, die täglich nicht einen Bissen essen noch einen Faden am Leibe tragen, den sie uns nicht gestohlen und geraubt haben durch ihren verdammten Wucher, leben also täglich von eitel Diebstahl und Raub mit Weib und Kind, als die Erzdiebe und Landräuber, in aller unbußfertigen Sicherheit. Denn ein Wucherer ist ein Erzdieb und Landräuber, der billig am Galgen siebenmal höher als andere Diebe hängen sollte.}
         {Même s’ils n’étaient pas si totalement aveugles, leur propre vie, et jusqu’à leur vie extérieure la plus grossière, devrait bien les convaincre de la pénitence qu’ils font. Car ils sont pleins de sorcellerie et de charlatanerie, de signes, de figures et du nom tétragramme (c’est-à-dire d’idolâtrie), pleins d’envie et d’orgueil, et avec cela ce ne sont que voleurs et brigands, qui chaque jour ne mangent pas une bouchée ni ne portent sur eux un fil qu’ils ne nous aient volé et ravi par leur maudite usure ; ils vivent ainsi chaque jour, avec femme et enfants, de purs vols et rapines, en archivoleurs et pilleurs du pays, dans une entière et impénitente sécurité. Car un usurier est un archivoleur et un pilleur du pays, qui devrait à bon droit être pendu au gibet sept fois plus haut que les autres voleurs.}
  \chunks{\textelp{}}
         {\textelp{}}
% Partie III : WA 53, p. 511-522
  \chunks{\subsection*{III. Die Lästerungen gegen Jesus und Maria}}
         {\subsection*{III. Les calomnies contre Jésus et Marie}}
% WA 53, p. 511-512
  \chunks{Wollen am Ende auch ihre Lügen wider die Personen sehen, welche doch die Lehre weder ärger noch besser machen, sie seien fromm oder böse. Nämlich was sie lügen von der Person unsers Herrn, item von seiner lieben Mutter und von uns und allen Christen. Das sind solche Lügen: wenn der Teufel der Lehre nichts kann anhaben, so legt er sich wider die Person, lügt, schmäht, flucht und tobt wider dieselbigen. Gleich wie der Papisten Beelzebub mir tat: da er meinem Evangelio nicht konnte widerstehen, schrieb er, ich hätte den Teufel, wäre ein Wechselbalg, meine liebe Mutter eine Hure und Bademagd; flugs, da er das hatte geschrieben, ward mein Evangelium zunichte und hatten die Papisten gewonnen. Also mußte Johannes der Täufer und Christus selbs den Teufel haben und Samariter heißen; flugs darauf ward Johannes’ und Christus’ Lehre unrecht und der Pharisäer Lehre recht. So ists allen Propheten gangen. Neulich itzt auch, da der Meuchelmordbrenner von Wolfenbüttel, der heiligen römischen Kirchen nach dem Kardinal zu Mainz einiges Heiligtum und Kleinod, die Person des Kurfürsten zu Sachsen und Landgrafen zu Hessen hatte geschmäht und geschändet aufs schändlichst, da waren sie flugs beide verloren; er aber, der heilige Mann, König über alle Könige, ward mit Rautenkranz und Gold gekrönet, so schwer, daß ers nicht ertragen kunnt und fliehen mußte.}
         {Nous voulons enfin voir aussi leurs mensonges contre les personnes, lesquelles pourtant ne rendent la doctrine ni pire ni meilleure, qu’elles soient pieuses ou méchantes : à savoir ce qu’ils mentent sur la personne de notre Seigneur, de même sur sa chère mère, et sur nous et tous les chrétiens. Ce sont des mensonges de cette sorte : quand le diable ne peut rien contre la doctrine, il s’en prend à la personne, il ment, injurie, maudit et fait rage contre elle. Ainsi m’a fait le Belzébuth des papistes : comme il ne pouvait résister à mon Évangile, il a écrit que j’avais le diable, que j’étais un changelin, que ma chère mère était une putain et une servante d’étuves ; aussitôt qu’il eut écrit cela, mon Évangile fut anéanti, et les papistes avaient gagné. Ainsi Jean-Baptiste et le Christ lui-même durent avoir le diable et être appelés Samaritains ; aussitôt la doctrine de Jean et du Christ fut fausse, et celle des pharisiens juste. C’est ce qui est arrivé à tous les prophètes. Tout récemment encore, quand le traître assassin et incendiaire de Wolfenbüttel\footnote{Henri~II, duc de Brunswick-Wolfenbüttel, prince catholique, adversaire de la ligue de Smalkalde ; Luther l’accusait d’avoir soudoyé des incendiaires. L’électeur de Saxe et le landgrave de Hesse l’ont chassé de son duché en 1542.}, l’unique sanctuaire et joyau de la sainte Église romaine après le cardinal de Mayence, eut injurié et outragé de la façon la plus ignominieuse la personne du prince électeur de Saxe et celle du landgrave de Hesse, aussitôt ils furent perdus tous les deux ; et lui, le saint homme, roi au-dessus de tous les rois, fut couronné de la couronne de rue et d’or\footnote{Le \foreignlanguage{german}{\emph{Rautenkranz}} (couronne de rue) figure sur les armoiries de la maison de Saxe : le duc n’a été \enquote{couronné} que par sa défaite.}, si lourde qu’il ne put la porter et dut s’enfuir.}
% WA 53, p. 512
  \chunks{Darum, wer eine böse Sache gewinnen will, der tue auch also, und wie die losen Zungendrescher vor Gericht tun, wenn sie die Silbersucht und das Guldenfieber bestehet: schelte und lüge getrost auf die Person, so ist die Sache gewonnen. Wie jene Mutter ihr Kind lehret: Lieber Sohn, kannst du nicht gewinnen, so trag Hader ein. Das heißen solche Lügen, da der Lügner nicht wähnet noch irret in der Hauptsachen (wie in Glaubenssachen auch geschicht), sondern selbs wohl weiß, daß er leugt und lügen will wider die Person, gedenkts auch nicht, weder mit Schein noch Wahrheit, zu beweisen, wie er auch nicht kann.}
         {C’est pourquoi, qui veut gagner une mauvaise cause n’a qu’à faire de même, comme font devant les tribunaux les bavards sans scrupules quand les prend le mal de l’argent et la fièvre des florins : qu’il injurie et mente hardiment contre la personne, et la cause est gagnée. Comme cette mère qui enseignait à son enfant : Mon fils, si tu ne peux gagner, sème la zizanie. Voilà ce que j’appelle des mensonges de cette sorte : le menteur ne se fait pas d’illusion et ne se trompe pas sur le fond (comme il arrive aussi en matière de foi), mais il sait fort bien lui-même qu’il ment et veut mentir contre la personne ; et il ne songe même pas à le prouver, ni en apparence ni en vérité, comme d’ailleurs il ne le peut.}
% WA 53, p. 512
  \chunks{Also tun die Juden in diesem Fall auch: getrost schelten sie, lügen und fluchen der Person, wider ihr eigen Gewissen; damit haben sie die Sachen längst gewonnen, daß sie Gott hat müssen erhören, sitzen zu Jerusalem in einer gülden Stadt, nun wohl 1500 Jahr, wie wir vor Augen sehen, sind Herren der Welt, laufen alle Heiden zu mit ihrem \emph{Hemdath}, Röcken, Bruchen und Schuhen, lassen sich würgen von den edlen Fürsten und Herren von Israel, geben ihnen Land und Leute und alles, was sie haben, wie sie den Gojim fluchen, speien und maledeien.}
         {Ainsi font aussi les Juifs en ce cas : ils injurient hardiment la personne, mentent sur elle et la maudissent, contre leur propre conscience ; par là ils ont depuis longtemps gagné leur cause, si bien que Dieu a dû les exaucer : ils siègent à Jérusalem dans une ville d’or, voilà bien 1500 ans, comme nous le voyons de nos yeux ; ils sont les maîtres du monde ; tous les gentils accourent à eux avec leur \emph{hemdath}\footnote{Jeu sur \emph{hemdath} (Aggée~2,~7) et \foreignlanguage{german}{\emph{Hemd}} (\enquote{chemise}), qui ouvre l’énumération des vêtements (glose de WA).}, leurs robes, leurs braies et leurs souliers, se laissent égorger par les nobles princes et seigneurs d’Israël, leur donnent leurs terres, leurs sujets et tout ce qu’ils ont, tandis qu’eux maudissent, conspuent et exècrent les goyim.}
% WA 53, p. 512
  \chunks{Und wo sie nicht so schändlich lögen, fluchten, speieten, lästerten, schmäheten die Personen, kannst du denken, so hätte sie Gott nicht erhöret, und hätten ihre Sache längst verloren, müßten nicht zu Jerusalem Herren sein, sondern in der Welt zerstreuet sein, Jerusalem nicht sehen, sich mit Lügen, Triegen, Stehlen, Rauben, Wuchern und allerlei Untugend unter den verfluchten Gojim nähren.}
         {Et s’ils ne mentaient pas si ignominieusement, s’ils ne maudissaient, ne conspuaient, ne blasphémaient et n’injuriaient pas ainsi les personnes, tu peux bien penser que Dieu ne les aurait pas exaucés : ils auraient depuis longtemps perdu leur cause, ne seraient pas les seigneurs de Jérusalem, mais seraient dispersés dans le monde, sans voir Jérusalem, et devraient vivre de mensonge, de tromperie, de vol, de rapine, d’usure et de toutes sortes de vices parmi les goyim maudits.}
  \chunks{Solch gewaltig Ding ists, wenn einer fluchen kann den Personen, so die Sache böse und verloren ist. Darum laß dir der Juden Exempel nicht ein schlecht Ding sein, wo du eine böse Sache hast. Sie sind die edlen Fürsten von Israel, die alles vermögen, wenn sie die Sache verloren haben, und den Gojim sehr fluchen können.}
         {Tant c’est chose puissante que de savoir maudire les personnes quand la cause est mauvaise et perdue ! Ne tiens donc pas pour peu de chose l’exemple des Juifs, quand tu as une mauvaise cause. Ce sont les nobles princes d’Israël, qui peuvent tout quand ils ont perdu leur cause, et qui savent si bien maudire les goyim.}
  \chunks{\textelp{}}
         {\textelp{}}
% WA 53, p. 515
  \chunks{Nun, was hat ihnen doch die arme Magd Maria getan? Womit wollen sie beweisen, daß sie eine Hure sei? Sie hat nichts getan, denn einen Sohn getragen, der Jesus heißt. Ists denn also ein böse Ding, daß ein junges Eheweib ein Kind gebiert? Oder sollen sie alle Huren heißen, die ein Kind tragen? Wo wollen ihre Weiber und sie selbs bleiben? Sind sie auch alle Huren und Hurenkinder? Ja, ihr verfluchten Gojim, das ist ein anders. Wisset ihr nicht, die Juden sind Abrahams edles Blut, beschnitten und Könige in Himmel und Erden; was sie sagen, das ist recht. Wenn nun unter den verfluchten Gojim eine Jungfrau wäre, so rein und heilig als der Engel Gabriel ist, und der geringsten edlen Fürsten einer spräche, sie wäre eine Erzhure und ärger denn der Teufel, so müßt es also sein, dürft keiner ander Beweisung denn dieser: Es hat ein edles Maul vom Geblüt Abraham gesagt, wer tar nein dazu sagen? Widerum, so eine Erzhure vom edlen Blut der Juden wäre, so häßlich als der Teufel selbs, wenns den edlen Herren gefiele, so ist sie reiner weder kein Engel. Denn die Macht stehet bei den edlen beschnittenen Herren, die verfluchten Gojim zu belügen, schmähen, schänden, lästern und fluchen, wie sie wollen, widerum sich selbs zu segnen, ehren, preisen und loben, wenns auch Gott anders wollte. Meinst du, es sei ein solch schlecht Ding um einen Juden? Gott im Himmel und alle Engel müssen lachen und tanzen, wenn sie einen Juden hören einen Furz lassen, auf daß du verfluchter Goi hinfort wissest, wie ein herrlich Ding es um einen Juden ist. Wie türsten sie sonst so kühnlich Maria eine Hure schelten, an der sie keine Schuld wissen, wo sie nicht solche Macht hätten, Gott und sein Gebot mit Füßen zu treten?}
         {Mais que leur a donc fait la pauvre servante Marie ? Avec quoi veulent-ils prouver qu’elle soit une putain ? Elle n’a rien fait, sinon porter un fils qui s’appelle Jésus. Est-ce donc une si mauvaise chose qu’une jeune épouse mette au monde un enfant ? Ou faut-il appeler putains toutes celles qui portent un enfant ? Que deviendront alors leurs femmes, et eux-mêmes ? Sont-ils donc tous, eux aussi, des putains et des enfants de putains ? Oui, mais, goyim maudits, c’est autre chose. Ne savez-vous pas que les Juifs sont le noble sang d’Abraham, circoncis, rois au ciel et sur la terre ? Ce qu’ils disent est juste. Si donc il y avait parmi les goyim maudits une vierge aussi pure et sainte que l’ange Gabriel, et que le moindre des nobles princes déclarât qu’elle est une archiputain, et pire que le diable, il faudrait qu’il en soit ainsi ; il ne serait besoin d’aucune autre preuve que celle-ci : une noble gueule du sang d’Abraham l’a dit, qui oserait dire non ? À l’inverse, s’il y avait une archiputain du noble sang des Juifs, aussi laide que le diable lui-même, il suffit qu’il plaise aux nobles seigneurs, et la voilà plus pure qu’aucun ange. Car le pouvoir appartient aux nobles seigneurs circoncis de mentir sur les goyim maudits, de les injurier, les outrager, les blasphémer et les maudire à leur gré, et, à l’inverse, de se bénir, de s’honorer, de se louer et de se glorifier eux-mêmes, quand bien même Dieu voudrait autre chose. Crois-tu qu’un Juif soit si peu de chose ? Dieu au ciel et tous les anges doivent rire et danser quand ils entendent un Juif lâcher un pet, afin que tu saches désormais, goy maudit, quelle chose magnifique est un Juif. Comment, sinon, oseraient-ils traiter si hardiment Marie de putain, alors qu’ils ne lui connaissent aucune faute, s’ils n’avaient le pouvoir de fouler aux pieds Dieu et son commandement ?}
  \chunks{\textelp{}}
         {\textelp{}}
% WA 53, p. 520
  \chunks{Nun siehe, welch eine feine, dicke, fette Lüge das ist, da sie klagen, sie seien bei uns gefangen. Es sind über 1400 Jahr, daß Jerusalem zerstöret ist, und wir Christen zu der Zeit schier 300 Jahr lang von den Juden gemartert und verfolget sind in aller Welt (wie droben gesagt), daß wir wohl möchten klagen, sie hätten uns Christen zu der Zeit gefangen und getötet, wie es die helle Wahrheit ist. Dazu wissen wir noch heutiges Tages nicht, welcher Teufel sie her in unser Land bracht hat; wir haben sie zu Jerusalem nicht geholet.}
         {Vois donc quel beau, gros et gras mensonge c’est, quand ils se plaignent d’être captifs chez nous. Il y a plus de 1400 ans que Jérusalem est détruite, et nous, chrétiens, avons été en ce temps-là, pendant près de 300 ans, martyrisés et persécutés par les Juifs dans le monde entier (comme on l’a dit plus haut), si bien que nous pourrions à bon droit nous plaindre qu’ils nous aient alors tenus captifs et mis à mort, ce qui est la vérité toute claire. De plus, nous ne savons toujours pas aujourd’hui quel diable les a amenés dans notre pays ; ce n’est pas nous qui sommes allés les chercher à Jérusalem.}
  \chunks{Zudem hält sie noch itzt niemand, Land und Straßen stehen ihnen offen, mögen ziehen in ihr Land, wenn sie wollen; wir wollten gern Geschenk dazu geben, daß wir ihr los wären. Denn sie sind uns eine schwere Last, wie eine Plage, Pestilenz und eitel Unglück in unserm Lande. Zum Wahrzeichen sind sie oft mit Gewalt vertrieben (schweige, daß wir sie sollten halten): aus Frankreich (das sie Zarpath nennen aus Obadja) als einem feinen sonderlichen Nest sind sie vertrieben; itzt neulich sind sie von dem lieben Kaiser Karolo aus Hispanien (welchs sie Sepharad auch aus Obadja nennen), dem allerbesten Nest, vertrieben; und dies Jahr aus der ganzen Böhmischen Krone, da sie doch zu Prag auch der besten Nester eines hatten; item aus Regensburg, Magdeburg und mehr Orten bei meinem Leben vertrieben.}
         {En outre, personne ne les retient aujourd’hui ; pays et routes leur sont ouverts, ils peuvent s’en aller dans leur pays quand ils veulent ; nous leur donnerions volontiers des présents par-dessus le marché pour être débarrassés d’eux. Car ils sont pour nous un lourd fardeau, comme une plaie, une peste et un pur malheur dans notre pays. La preuve en est qu’on les a souvent chassés par la force, bien loin qu’on ait voulu les retenir : de France (qu’ils appellent Zarpath, d’après Abdias\footnote{Abdias~20 nomme Sarepta et Sepharad ; la tradition juive médiévale en a fait les noms de la France (\emph{Tsarfat}) et de l’Espagne (\emph{Sefarad}).}), qui était pour eux un beau nid de choix, ils ont été chassés ; tout récemment, notre cher empereur Charles\footnote{Erreur de Luther : l’expulsion des Juifs d’Espagne date de 1492, sous Ferdinand et Isabelle ; Charles Quint ne régna sur l’Espagne qu’à partir de 1516.} les a chassés d’Espagne (qu’ils appellent Sepharad, également d’après Abdias), le meilleur de tous les nids ; et cette année, de toute la couronne de Bohême, alors qu’ils avaient pourtant à Prague l’un des meilleurs nids ; de même de Ratisbonne, de Magdebourg et d’autres lieux, de mon vivant.}
% WA 53, p. 521
  \chunks{Heißt das gefangen halten, wenn man einen nicht leiden kann im Lande oder Hause? Jawohl, sie halten uns Christen in unserm eigen Lande gefangen: sie lassen uns arbeiten im Nasenschweiß, Geld und Gut gewinnen, sitzen sie dieweil hinter dem Ofen, faulenzen, pompen und braten Birn, fressen, saufen, leben sanft und wohl von unserm erarbeiteten Gut, haben uns und unser Güter gefangen durch ihren verfluchten Wucher, spotten dazu und speien uns an, daß wir arbeiten und sie faule Junker lassen sein von dem Unsern und in dem Unsern; sind also unsere Herren, wir ihre Knechte mit unserm eigen Gut, Schweiß und Arbeit, fluchen darnach unserm Herrn und uns zu Lohn und zu Dank. Sollt der Teufel hie nicht lachen und tanzen, wenn er solch fein Paradies bei uns Christen haben kann, daß er durch die Juden, seine Heiligen, das Unser frisset und uns zu Lohn Maul und Nasen voll tut, spottet und flucht Gott und Menschen dazu?}
         {Est-ce là tenir captif, quand on ne peut souffrir quelqu’un dans son pays ou dans sa maison ? Oui, assurément, ils nous tiennent captifs, nous chrétiens, dans notre propre pays : ils nous laissent travailler à la sueur de notre front\footnote{\foreignlanguage{german}{\emph{Im Nasenschweiß}}, littéralement \enquote{à la sueur du nez} : variante populaire de la formule de Genèse~3,~19 (\enquote{à la sueur de ton visage}), que Luther applique ici aux chrétiens.}, gagner argent et biens, tandis qu’ils sont assis derrière le poêle, paressent, pètent et font rôtir des poires, bâfrent, s’enivrent, vivent doucement et bien de nos biens durement gagnés ; ils nous ont faits captifs, nous et nos biens, par leur maudite usure ; ils se moquent de nous par-dessus le marché et nous conspuent, parce que nous travaillons et les laissons faire les hobereaux fainéants, de notre bien et dans notre bien. Ils sont donc nos maîtres, et nous leurs valets, avec notre propre bien, notre sueur et notre travail ; et ensuite ils maudissent notre Seigneur et nous-mêmes, pour salaire et pour remerciement. Le diable n’aurait-il pas de quoi rire et danser, quand il peut avoir chez nous, chrétiens, un si beau paradis : par les Juifs, ses saints, il dévore notre bien, et pour salaire il nous en met plein la gueule, se moque de Dieu et des hommes et les maudit par-dessus le marché ?}
  \chunks{\textelp{}}
         {\textelp{}}
% WA 53, p. 522
  \chunks{Aus diesem allen sehen wir Christen (denn sie, die Juden, könnens nicht sehen), welch ein schrecklicher Zorn Gottes über dies Volk gangen und ohn Aufhören gehet, welch ein Feuer und Glut brennet da, und was die gewinnen, so Christum und seinen Christen fluchen oder feind sind. O lieben Christen, laßt uns solch greulich Exempel zu Herzen nehmen, wie S. Paulus Röm. 11 sagt, und Gott fürchten, daß wir nicht auch zuletzt in solchen oder noch ärgern Zorn fallen, sondern (wie wir droben auch gesagt) sein göttlich Wort ehren und die Zeit der Gnaden nicht versäumen, wie es bereits der Mahomet und Papst versäumet haben und nicht viel besser denn die Juden worden sind.}
         {Par tout cela, nous voyons, nous chrétiens (car eux, les Juifs, ne peuvent le voir), quelle terrible colère de Dieu s’est abattue sur ce peuple et s’abat sans relâche, quel feu, quelle braise brûle là, et ce que gagnent ceux qui maudissent le Christ et ses chrétiens ou leur sont hostiles. Ô chers chrétiens, prenons à cœur un exemple si effroyable, comme le dit saint Paul en Romains~11, et craignons Dieu, de peur que nous aussi ne tombions à la fin dans une semblable colère, ou pire encore ; mais (comme nous l’avons dit aussi plus haut) honorons sa parole divine et ne laissons pas passer le temps de la grâce, comme l’ont déjà laissé passer Mahomet et le pape, qui ne sont pas devenus beaucoup meilleurs que les Juifs.}
% Partie IV : WA 53, p. 522-541
  \chunks{\subsection*{IV. Was mit den Juden zu tun ist}}
         {\subsection*{IV. Ce qu’il faut faire des Juifs}}
% WA 53, p. 522
  \chunks{Was wollen wir Christen nun tun mit diesem verworfenen, verdammten Volk der Juden? Zu leiden ists uns nicht, nachdem sie bei uns sind und wir solch Lügen, Lästern und Fluchen von ihnen wissen, damit wir uns nicht teilhaftig machen aller ihrer Lügen, Flüche und Lästerung. So können wir das unlöschliche Feuer göttlichs Zorns (wie die Propheten sagen) nicht löschen noch die Juden bekehren. Wir müssen mit Gebet und Gottesfurcht eine scharfe Barmherzigkeit üben, ob wir doch etliche aus der Flamme und Glut erretten könnten. Rächen dürfen wir uns nicht; sie haben die Rache am Halse, tausendmal ärger, denn wir ihnen wünschen mögen.}
         {Que voulons-nous donc faire, nous chrétiens, de ce peuple rejeté et damné des Juifs ? Il ne nous est pas possible de le tolérer, puisqu’ils sont chez nous et que nous savons de leur part de tels mensonges, blasphèmes et malédictions, sous peine de nous rendre complices de tous leurs mensonges, malédictions et blasphèmes. Or nous ne pouvons ni éteindre le feu inextinguible de la colère divine (comme disent les prophètes), ni convertir les Juifs. Il nous faut, avec prière et crainte de Dieu, exercer une âpre miséricorde, pour voir si nous ne pourrions pas du moins en sauver quelques-uns de la flamme et de la braise. Nous venger, nous n’en avons pas le droit : la vengeance leur pend au nez, mille fois pire que nous ne pourrions la leur souhaiter.}
% WA 53, p. 522
  \chunks{Ich will meinen treuen Rat geben.}
         {Je veux donner mon conseil loyal.}
% WA 53, p. 523
  \chunks{Erstlich, daß man ihre Synagoga oder Schule mit Feuer anstecke und, was nicht verbrennen will, mit Erden überhäufe und beschütte, daß kein Mensch einen Stein oder Schlacke davon sehe ewiglich. Und solchs soll man tun unserm Herrn und der Christenheit zu Ehren, damit Gott sehe, daß wir Christen seien und solch öffentlich Lügen, Fluchen und Lästern seines Sohnes und seiner Christen wissentlich nicht geduldet noch gewilliget haben.}
         {Premièrement, qu’on mette le feu à leurs synagogues, et que ce qui ne veut pas brûler, on le recouvre et l’enfouisse sous la terre, afin que personne n’en voie plus jamais une pierre ni une scorie. Et cela, on doit le faire en l’honneur de notre Seigneur et de la chrétienté, afin que Dieu voie que nous sommes chrétiens et que nous n’avons pas sciemment toléré ni approuvé de tels mensonges, malédictions et blasphèmes publics contre son Fils et contre ses chrétiens.}
% WA 53, p. 523
  \chunks{Denn was wir bisher aus Unwissenheit geduldet (ich habs selbs nicht gewußt), wird uns Gott verzeihen. Nun wirs aber wissen und sollten darüber, frei vor unser Nasen, den Juden ein solch Haus schützen und schirmen, darin sie Christum und uns belügen, lästern, fluchen, anspeien und schänden (wie droben gehöret), das wäre eben so viel, als täten wirs selbs, und viel ärger, wie man wohl weiß.}
         {Car ce que nous avons toléré jusqu’ici par ignorance (je ne le savais pas moi-même), Dieu nous le pardonnera. Mais maintenant que nous le savons, si nous devions malgré cela, ouvertement et sous notre nez, protéger et abriter pour les Juifs une telle maison, où ils mentent sur le Christ et sur nous, nous blasphèment, nous maudissent, nous conspuent et nous outragent (comme on l’a entendu plus haut), ce serait exactement comme si nous le faisions nous-mêmes, et bien pire, comme on le sait bien.}
% WA 53, p. 523
  \chunks{Moses schreibt Deut. 13, daß, wo eine Stadt Abgötterei triebe, sollt man sie mit Feuer ganz verstören und nichts davon behalten. Und wenn er itzt lebete, so würde er der erste sein, der die Judenschulen und -häuser ansteckte. Denn er hat gar hart geboten Deut. 4 und 12, sie sollen nichts zu- noch abtun von seinem Gesetze. Und Samuel sagt 1. Sam. 15, es sei Abgötterei, Gott nicht gehorchen. Nun ist der Juden Lehre itzt nichts anders denn eitel Zusätze der Rabbinen und Abgötterei des Ungehorsams, daß Mose ganz unkenntlich bei ihnen worden ist (wie gesagt), gleich wie bei uns unter dem Papsttum die Biblia unkenntlich worden ist. Daß also auch Mosi halben ihre Schulen nicht zu leiden sind, den sie eben so wohl schänden als uns, und nicht not ist, daß sie zu solcher Abgötterei eigene, freie Kirchen haben sollten.}
         {Moïse écrit, en Deutéronome~13, que si une ville pratiquait l’idolâtrie, on devait la détruire entièrement par le feu et n’en rien garder. Et s’il vivait aujourd’hui, il serait le premier à mettre le feu aux synagogues et aux maisons des Juifs. Car il a très sévèrement ordonné, en Deutéronome~4 et~12, qu’on n’ajoute ni ne retranche rien à sa loi. Et Samuel dit, en 1~Samuel~15, que ne pas obéir à Dieu est idolâtrie. Or la doctrine des Juifs n’est plus aujourd’hui que pures additions des rabbins et idolâtrie de la désobéissance, si bien que Moïse leur est devenu entièrement inconnu (comme on l’a dit), de même que chez nous, sous la papauté, la Bible était devenue inconnue. Ainsi, à cause de Moïse aussi, il ne faut pas tolérer leurs synagogues, puisqu’ils l’outragent tout autant que nous ; et il n’est pas nécessaire qu’ils aient, pour une telle idolâtrie, des églises à eux, en toute liberté.}
% WA 53, p. 523
  \chunks{Zum andern, daß man auch ihre Häuser desgleichen zerbreche und zerstöre. Denn sie treiben eben dasselbige drinnen, das sie in ihren Schulen treiben. Dafür mag man sie etwa unter ein Dach oder Stall tun, wie die Zigeuner, auf daß sie wissen, sie seien nicht Herren in unserm Lande, wie sie rühmen, sondern im Elend und gefangen, wie sie ohn Unterlaß vor Gott über uns Zeter schreien und klagen.}
         {Deuxièmement, qu’on démolisse et détruise aussi, de même, leurs maisons. Car ils y pratiquent exactement la même chose que dans leurs synagogues. À la place, on pourra les mettre par exemple sous un toit ou dans une étable, comme les Tsiganes, afin qu’ils sachent qu’ils ne sont pas les maîtres dans notre pays, comme ils s’en vantent, mais en exil et captifs, comme ils ne cessent de le déplorer devant Dieu en criant haro sur nous.}
% WA 53, p. 523
  \chunks{Zum dritten, daß man ihnen nehme alle ihre Betbüchlein und Talmudisten, darin solche Abgötterei, Lügen, Fluch und Lästerung gelehret wird.}
         {Troisièmement, qu’on leur prenne tous leurs livrets de prières et leurs Talmuds, où l’on enseigne une telle idolâtrie, de tels mensonges, malédictions et blasphèmes.}
% WA 53, p. 523-524
  \chunks{Zum vierten, daß man ihren Rabbinen bei Leib und Leben verbiete, hinfort zu lehren. Denn solch Amt haben sie mit allem Recht verloren, weil sie die armen Juden mit dem Spruch Mosi, Deut. 17, gefangen halten, da er gebeut, sie sollen ihren Lehrern gehorchen bei Verlust Leibs und Seelen, so doch Mose klärlich daselbst beisetzt: \enquote{Was sie dich lehren nach dem Gesetz des HERRN}. Solchs übergehen die Bösewichter und brauchen des armen Volks Gehorsam zu ihrem Mutwillen wider das Gesetz des HERRN, gießen ihnen solche Gift, Fluch und Lästerung ein. Gleich wie uns der Papst mit dem Spruch Matth. 16: \enquote{Du bist Petrus etc.} gefangen hielt, daß wir alles mußten glauben, was er uns für Log und Trog aus seinem Teufelskopf und nicht nach Gottes Wort uns lehret, darüber er das Amt zu lehren verloren hat.}
         {Quatrièmement, qu’on interdise à leurs rabbins, sous peine de la vie, d’enseigner désormais. Car ils ont perdu cette charge à très bon droit, puisqu’ils tiennent captifs les pauvres Juifs par la parole de Moïse, Deutéronome~17, où il ordonne d’obéir à leurs maîtres sous peine de perdre corps et âme, alors que Moïse ajoute clairement au même endroit : \enquote{ce qu’ils t’enseignent selon la loi du SEIGNEUR}. Cela, les scélérats le passent sous silence ; ils usent de l’obéissance du pauvre peuple pour leur bon plaisir, contre la loi du SEIGNEUR, et lui versent un tel poison, de telles malédictions et de tels blasphèmes. De même le pape nous tenait-il captifs par la parole de Matthieu~16 : \enquote{Tu es Pierre, etc.}, si bien que nous devions croire tout ce qu’il nous enseignait de menteries et de tromperies, sorties de sa tête de diable, et non selon la parole de Dieu ; c’est pourquoi il a perdu la charge d’enseigner.}
  \chunks{\textelp{}}
         {\textelp{}}
% WA 53, p. 524
  \chunks{Zum fünften, daß man den Juden das Geleit und Straße ganz und gar aufhebe. Denn sie haben nichts auf dem Lande zu schaffen, weil sie nicht Herren noch Amtleute noch Händler oder desgleichen sind; sie sollen daheime bleiben. Ich lasse mir sagen, es solle ein reicher Jude itzt auf dem Lande reiten mit zwölf Pferden (der will ein Kochab werden) und wuchert Fürsten, Herren, Land und Leute aus, daß große Herren scheel dazu sehen. Werdet ihr Fürsten und Herren solchen Wucherern nicht die Straße legen ordentlicher Weise, so möcht sich etwa eine Reuterei sammeln wider sie, weil sie aus diesem Büchlein lernen werden, was die Juden sind und wie man mit ihnen umgehen und ihr Wesen nicht schützen solle. Denn ihr sollt und könnt sie auch nicht schützen, ihr wollt denn vor Gott alles ihres Greuels teilhaftig sein. Was daraus Guts kommen möchte, das wollet wohl bedenken und verkommen.}
         {Cinquièmement, qu’on retire entièrement aux Juifs le sauf-conduit et l’usage des routes. Car ils n’ont rien à faire dans le pays, puisqu’ils ne sont ni seigneurs, ni officiers, ni marchands, ni rien de tel ; qu’ils restent chez eux. J’entends dire qu’un riche Juif chevauche maintenant à travers le pays avec douze chevaux (celui-là veut devenir un Kochab\footnote{\emph{Kochab}, \enquote{étoile} : Bar Kokhba, chef de la révolte juive de 132-135 contre Rome, salué comme Messie par Rabbi Akiba ; Luther en a parlé plus haut dans le traité.}) et ruine par l’usure princes, seigneurs, pays et gens, si bien que de grands seigneurs le voient d’un mauvais œil. Si vous, princes et seigneurs, ne barrez pas la route à de tels usuriers par les voies régulières, il pourrait bien se former un jour contre eux une troupe de cavaliers, car ils apprendront dans ce petit livre ce que sont les Juifs, comment il faut agir avec eux, et qu’il ne faut pas protéger leurs agissements. Car vous ne devez ni ne pouvez les protéger, à moins de vouloir, devant Dieu, prendre part à toutes leurs abominations. Le bien qui pourrait sortir de là, pesez-le bien, et prévenez-le.}
% WA 53, p. 524
  \chunks{Zum sechsten, daß man ihnen den Wucher verbiete und nehme ihnen alle Barschaft und Kleinod an Silber und Gold und lege es beiseit zu verwahren. Und ist dies die Ursache: alles, was sie haben (wie droben gesagt), haben sie uns gestohlen und geraubt durch ihren Wucher, weil sie sonst keine andere Nahrung haben. Solch Geld sollt man dazu brauchen (und nicht anders), wo ein Jude sich ernstlich bekehret, daß man ihm davon vor die Hand gebe hundert, zwei-, dreihundert Gulden nach Gelegenheit der Person, damit er eine Nahrung für sein arm Weib und Kindlein anfahen möge und die Alten oder Gebrechlichen damit unterhalte. Denn solch böse gewonnen Gut ist verflucht, wo mans nicht mit Gottes Segen in guten, nötigen Brauch wendet.}
         {Sixièmement, qu’on leur interdise l’usure, qu’on leur prenne tout leur argent comptant et leurs joyaux d’argent et d’or, et qu’on les mette de côté en dépôt. En voici la raison : tout ce qu’ils ont (comme on l’a dit plus haut), ils nous l’ont volé et ravi par leur usure, puisqu’ils n’ont pas d’autre gagne-pain. Cet argent, on devrait l’employer à ceci, et à rien d’autre : quand un Juif se convertit sérieusement, qu’on lui en remette en main cent, deux cents, trois cents florins, selon la condition de la personne, pour qu’il puisse commencer un métier pour sa pauvre femme et ses petits enfants, et entretenir avec cela les vieux ou les infirmes. Car un bien si mal acquis est maudit, si on ne le tourne pas, avec la bénédiction de Dieu, à un usage bon et nécessaire.}
  \chunks{\textelp{}}
         {\textelp{}}
% WA 53, p. 525-526
  \chunks{Zum siebenten, daß man den jungen, starken Juden und Jüdin in die Hand gebe Flegel, Axt, Karst, Spaten, Rocken, Spindel und lasse sie ihr Brot verdienen im Schweiß der Nasen, wie Adams Kindern auferlegt ist, Gen. 3. Denn es taugt nicht, daß sie uns verfluchten Gojim wollten lassen im Schweiß unsers Angesichts arbeiten, und sie, die heiligen Leute, woltens hinter dem Ofen mit faulen Tagen, feisten und pompen verzehren und darauf rühmen lästerlich, daß sie der Christen Herren wären von unserm Schweiß, sondern man müßte ihnen das faule Schelmenbein aus dem Rücken vertreiben.}
         {Septièmement, qu’on mette entre les mains des jeunes Juifs et Juives vigoureux le fléau, la hache, la houe, la bêche, la quenouille, le fuseau, et qu’on les laisse gagner leur pain à la sueur de leur front\footnote{\foreignlanguage{german}{\emph{Im Schweiß der Nasen}} : même expression qu’à la p.~521, variante populaire de Genèse~3,~19 ; Luther renvoie lui-même à Genèse~3 et dit juste après, pour les chrétiens, \foreignlanguage{german}{\emph{im Schweiß unsers Angesichts}} (\enquote{à la sueur de notre visage}).}, comme il a été imposé aux enfants d’Adam, Genèse~3. Car il ne convient pas qu’ils veuillent nous laisser travailler, nous goyim maudits, à la sueur de notre front, tandis qu’eux, les saints, consommeraient le tout derrière le poêle, en journées de paresse, à vesser et à péter, et se vanteraient par-dessus le marché, en blasphémateurs, d’être les maîtres des chrétiens grâce à notre sueur ; il faudrait au contraire leur chasser de l’échine leurs os paresseux de vauriens.}
% WA 53, p. 526
  \chunks{Besorgen wir uns aber, daß sie uns möchten an Leib, Weib, Kind, Gesind, Vieh etc. Schaden tun, wenn sie uns dienen oder arbeiten sollten, weil es wohl zu vermuten ist, daß solch edle Herren der Welt und giftige, bittere Würme, keiner Arbeit gewohnet, gar ungern sich so hoch demütigen würden unter die verfluchten Gojim, so laßt uns bleiben bei gemeiner Klugheit der andern Nation, als Frankreich, Hispanien, Böhmen etc., und mit ihnen rechnen, was sie uns abgewuchert, und darnach gütlich geteilet, sie aber immer zum Land ausgetrieben. Denn, wie gehört, Gottes Zorn ist so groß über sie, daß sie durch sanfte Barmherzigkeit nur ärger und ärger, durch scharfe aber wenig besser werden. Drum immer weg mit ihnen.}
         {Mais si nous craignons qu’ils puissent nous nuire, à nous, à nos femmes, nos enfants, nos domestiques, notre bétail, etc., s’ils devaient nous servir ou travailler pour nous, car il est bien à présumer que de si nobles seigneurs du monde, vers venimeux et amers, qui ne sont habitués à aucun travail, s’humilieraient de très mauvais gré à ce point sous les goyim maudits, alors restons-en à la sagesse commune des autres nations, comme la France, l’Espagne, la Bohême, etc. : faisons le compte avec eux de ce qu’ils nous ont extorqué par l’usure, puis partageons à l’amiable, mais chassons-les du pays une fois pour toutes. Car, comme on l’a entendu, la colère de Dieu est si grande sur eux que la douce miséricorde ne les rend que pires et pires, et l’âpre guère meilleurs. Donc, en tout cas, dehors avec eux\footnote{Formule reprise, sous la forme \foreignlanguage{german}{\emph{Weg mit ihnen!}}, dans le titre de la brochure de Martin Sasse, \foreignlanguage{german}{\emph{Martin Luther über die Juden: Weg mit ihnen!}} (1938), et, sous la forme \foreignlanguage{german}{\emph{Drum immer weg mit ihnen!}}, comme titre d’une exposition récente sur Luther et les Juifs.} !}
% WA 53, p. 526
  \chunks{Ich höre sagen, daß die Juden große Summa Geldes geben und damit den Herrschaften nütze sind. Ja, wovon geben sie es? Nicht von dem Ihren, sondern von der Herrschaft und Untertanen Güter, welche sie durch Wucher stehlen und rauben. Und nehmen also die Herrschaften von ihren Untertanen, was die Juden geben, das ist: die Untertanen müssen Geld zugeben und sich schinden lassen für die Juden, damit sie im Lande bleiben, getrost und frei lügen, lästern, fluchen und stehlen können.}
         {J’entends dire que les Juifs donnent de grosses sommes d’argent et sont ainsi utiles aux seigneurs. Oui, mais avec quoi les donnent-ils ? Non avec leur bien, mais avec les biens des seigneurs et des sujets, qu’ils volent et pillent par l’usure. Ainsi les seigneurs prennent à leurs sujets ce que donnent les Juifs ; c’est-à-dire que les sujets doivent payer en plus et se laisser écorcher pour les Juifs, afin que ceux-ci puissent rester dans le pays, et mentir, blasphémer, maudire et voler en toute tranquillité et en toute liberté.}
  \chunks{Sollten die verzweifelten Juden des nicht in die Faust lachen, daß wir uns so schändlich äffen und närren lassen und unser Geld geben, daß sie im Lande bleiben und alle Bosheit treiben mögen, über das noch reich dazu werden von unserm Schweiß und Blut, wir aber arm und von ihnen ausgesogen werden? Wenn das recht ist, daß ein Knecht, ja ein Gast möge seinem Herrn oder Wirt jährlich zehn Gulden geben und dafür tausend stehlen, so ist der Knecht und Gast leicht und bald reich, der Herr und Wirt in kurz ein Bettler worden.}
         {Ces scélérats de Juifs n’auraient-ils pas de quoi rire sous cape, de nous voir nous laisser si honteusement berner et duper, et donner notre argent pour qu’ils restent dans le pays et puissent commettre toute méchanceté, et que par-dessus le marché ils s’enrichissent de notre sueur et de notre sang, tandis que nous devenons pauvres, sucés par eux ? S’il est juste qu’un valet, voire un pensionnaire, puisse donner chaque année dix florins à son maître ou à son logeur, et lui en voler mille, alors le valet et le pensionnaire sont vite et aisément riches, et le maître et le logeur, en peu de temps, réduits à la mendicité.}
% WA 53, p. 526-527
  \chunks{Und wenn gleich die Juden von ihrem Eigen der Herrschaft solche Summa geben könnten, als nicht möglich ist, und sie damit uns abkaufen sollten Schutz und Schirm, öffentlich, frei in ihren Schulen unsern Herrn Christum so schändlich zu belügen, zu lästern, verspeien, verfluchen, dazu uns auch alles Unglück, daß wir alle erstochen werden und umkommen mit unserm Haman, Kaiser, Fürsten, Herren, Weib und Kindern, zu wünschen, das hieße wahrlich Christum, unsern Herrn, die ganze Christenheit samt dem ganzen Kaisertum, uns mit Weib und Kindern schändlich wohlfeil verkauft. Wie gar ein größer Heiliger würde hie der Verräter Judas gegen uns geschätzt werden? Ja, wenn ein jeglicher Jude (so viel ihr ist) jährlich hunderttausend Gulden geben könnte, so sollten wir doch nicht dafür gestatten, daß sie einen einigen Christen so frei zu lästern, zu fluchen, zu verspeien, auszuwuchern sollten Macht haben; es wäre noch viel zu wohlfeil verkauft. Wie viel unträglicher ists, daß wir den ganzen Christum und uns alle sollten mit unserm eigen Gelde kaufen lassen, den Juden zu lästern und zu fluchen, und sie zu Lohn davon noch reich und unser Junkherren machen, die uns dazu noch verlachten und sich in ihrem Mutwillen kitzelten. Das möchte dem Teufel und seinen Engeln ein recht Freudenspiel sein, des sie durch die Nasen lachen könnten, wie eine Sau ihre Ferkel anlacht, aber vor Gott einen rechten Zorn verdienen.}
         {Et quand bien même les Juifs pourraient donner aux seigneurs de telles sommes sur leur propre bien, ce qui n’est pas possible, et nous acheter par là protection et sauvegarde pour pouvoir, ouvertement et librement, dans leurs synagogues, mentir si ignominieusement sur notre Seigneur le Christ, le blasphémer, le conspuer, le maudire, et nous souhaiter en outre tous les malheurs, que nous soyons tous transpercés et que nous périssions avec notre Haman\footnote{Haman, l’ennemi des Juifs dans le livre d’Esther ; Luther prête aux Juifs l’usage d’appliquer ce nom aux souverains chrétiens (cf. p.~433).}, empereur, princes, seigneurs, femmes et enfants, ce serait vraiment avoir vendu honteusement, à vil prix, le Christ notre Seigneur, toute la chrétienté avec tout l’Empire, et nous-mêmes avec femmes et enfants. Combien le traître Judas passerait-il ici pour un plus grand saint que nous ! Oui, quand chaque Juif, autant qu’il y en a, pourrait donner chaque année cent mille florins, nous ne devrions pas pour autant leur permettre d’avoir le pouvoir de blasphémer, de maudire, de conspuer et de ruiner par l’usure un seul chrétien aussi librement : ce serait encore le vendre beaucoup trop bon marché. Combien plus insupportable est-il que nous devions nous laisser acheter, le Christ tout entier et nous tous, avec notre propre argent, pour être blasphémés et maudits par les Juifs, et qu’en récompense nous les enrichissions encore et en fassions nos hobereaux, qui par-dessus le marché se moqueraient de nous et se délecteraient de leur malice. Ce serait pour le diable et ses anges un vrai spectacle de réjouissance, dont ils pourraient ricaner du nez comme une truie sourit à ses porcelets ; mais ce serait mériter devant Dieu une juste colère.}
% WA 53, p. 527
  \chunks{Summa, lieben Fürsten und Herren, so Juden unter sich haben: ist euch solcher mein Rat nicht eben, so trefft einen bessern, daß ihr und wir alle der unleidlichen, teuflischen Last der Juden entladen werden und nicht vor Gott schuldig und teilhaftig werden alle der Lügen, des Lästerns, Speiens, Fluchens, so die rasenden Juden wider die Person unsers Herrn Jesu Christi, seiner lieben Mutter, aller Christen, aller Obrigkeit und unser selbs so frei und mutwilliglich treiben: keinen Schutz noch Schirm noch Geleit noch Gemeinschaft sie haben lassen, auch nicht euer und euer Untertanen Geld und Güter durch den Wucher ihnen dazu dienen und helfen lassen.}
         {En somme, chers princes et seigneurs qui avez des Juifs chez vous : si ce conseil que je vous donne ne vous convient pas, trouvez-en un meilleur, afin que vous et nous tous soyons déchargés du fardeau insupportable et diabolique des Juifs, et que nous ne devenions pas, devant Dieu, coupables et complices de tous les mensonges, blasphèmes, invectives et malédictions que les Juifs enragés profèrent si librement et si délibérément contre la personne de notre Seigneur Jésus-Christ, de sa chère mère, de tous les chrétiens, de toute autorité et de nous-mêmes. Ne leur laissez ni protection, ni sauvegarde, ni sauf-conduit, ni communauté, et ne laissez pas non plus votre argent et vos biens, ni ceux de vos sujets, les servir et les y aider par l’usure.}
% WA 53, p. 527
  \chunks{Wir haben zuvor eigener Sünde genug auf uns, noch vom Papsttum her, tun täglich viel dazu mit allerlei Undankbarkeit und Verachtung seines Worts und aller seiner Gnaden, daß nicht not ist, auch diese fremden, schändlichen Laster der Juden auf uns zu laden und ihnen dennoch Geld und Gut zu geben. Laßt uns denken, daß wir nun täglich wider den Türken streiten, da wir wohl Leichterung unser eigen Sünde und Besserung unsers Lebens zu dürfen. Ich will hiemit mein Gewissen gereinigt und entschuldigt haben, als der ichs treulich hab angezeigt und gewarnet.}
         {Nous avons déjà bien assez de nos propres péchés sur nous, qui nous restent de la papauté, et nous en ajoutons beaucoup chaque jour par toute sorte d’ingratitude et de mépris envers sa parole et toutes ses grâces, pour qu’il ne soit pas nécessaire de charger encore sur nous ces vices étrangers et ignominieux des Juifs et de leur donner en plus argent et biens. Songeons que nous combattons maintenant tous les jours contre le Turc\footnote{Allusion aux campagnes contre les Ottomans en Hongrie (prise de Buda, 1541). Pour Luther, le Turc est un fléau envoyé par Dieu, qu’on combat d’abord par la pénitence.}, et que nous aurions grand besoin d’être allégés de nos propres péchés et d’amender notre vie. Par là, je veux avoir purifié et déchargé ma conscience, pour avoir loyalement signalé et averti.}
% WA 53, p. 527-528
  \chunks{Und euch, meine lieben Herren und Freunde, so Pfarrherren und Prediger sind, will ich ganz treulich eures Amts hiemit erinnert haben, daß auch ihr euer Pfarrleute warnet für ihrem ewigen Schaden, wie ihr wohl zu tun wisset, nämlich daß sie sich vor den Juden hüten und sie meiden, wo sie können; nicht daß sie ihnen viel fluchen oder persönlich Leid tun sollten. Denn sie haben sich selbs allzu hoch verflucht und beleidigt, wenn sie den Mann Jesum von Nazareth, Marien Sohn, verfluchen, wie sie leider tun nun über 1400 Jahr. Die Obrigkeit lasse man hie mit ihnen gebaren, wie ich itzt gesagt. Es tu aber die Obrigkeit dazu oder nicht, daß dennoch ein jeder für sich selbs seines Gewissens wahrnehme und mache ihm eine solche definitio oder prosopopoeia eines Juden.}
         {Et vous, mes chers seigneurs et amis qui êtes pasteurs et prédicateurs, je veux par là vous avoir très loyalement rappelé votre charge : avertissez vous aussi vos paroissiens de leur dommage éternel, comme vous savez bien le faire, à savoir qu’ils se gardent des Juifs et les évitent où ils peuvent ; non qu’ils doivent beaucoup les maudire ou leur faire du mal en personne. Car ils se sont eux-mêmes maudits et outragés bien assez gravement, en maudissant l’homme Jésus de Nazareth, fils de Marie, comme ils le font hélas depuis plus de 1400 ans. Qu’on laisse ici l’autorité agir avec eux comme je viens de le dire. Mais que l’autorité s’en mêle ou non, que du moins chacun prenne garde pour lui-même à sa conscience, et se fasse d’un Juif une définition ou une prosopopée comme celle-ci\footnote{Termes de rhétorique latine : la \emph{definitio} définit une chose ; la \emph{prosopopoeia} fait parler une personne absente ou imaginaire. Luther annonce le discours que chacun doit se tenir à la vue d’un Juif.}.}
% WA 53, p. 528
  \chunks{Wenn du siehest oder denkest an einen Juden, so sprich bei dir selbs also: Siehe, das Maul, das ich da sehe, hat alle Sonnabend meinen lieben Herrn Jesum Christ, der mich mit seinem teuren Blut erlöset hat, verflucht und vermaledeiet und verspeiet, dazu gebetet und geflucht vor Gott, daß ich, mein Weib und Kind und alle Christen erstochen und aufs jämmerlichst untergangen wären, wollts selber gern tun, wo er könnte, daß er unser Güter besitzen möchte; hat auch vielleicht heute dieses Tages vielmal auf die Erden gespeiet über dem Namen Jesu (wie sie pflegen), daß ihm der Speichel noch im Maul und Bart hänget, wo er Raum hätte zu speien. Und ich sollte mit solchem verteufelten Maul essen, trinken oder reden? So möcht ich aus der Schüssel oder Kannen mich voller Teufel fressen und saufen, als der ich mich gewiß damit teilhaftig machte aller Teufel, so in den Juden wohnen und das teure Blut Christi verspeien. Da behüt mich Gott für.}
         {Quand tu vois un Juif ou que tu penses à un Juif, dis-toi en toi-même : Vois, cette gueule que je vois là a, chaque samedi, maudit, exécré et conspué mon cher Seigneur Jésus-Christ, qui m’a racheté de son sang précieux ; elle a en outre prié et maudit devant Dieu pour que moi, ma femme et mes enfants, et tous les chrétiens, nous soyons transpercés et périssions de la manière la plus misérable ; il le ferait volontiers lui-même, s’il le pouvait, pour posséder nos biens ; peut-être a-t-il encore aujourd’hui même craché bien des fois à terre sur le nom de Jésus (comme ils en ont coutume\footnote{Allusion probable à la coutume de cracher à un passage de la prière \emph{Aleinu}, que des polémistes chrétiens présentaient comme un crachat sur le nom de Jésus.}), si bien que la salive lui pend encore à la gueule et à la barbe, pour peu qu’il en ait eu l’occasion. Et je devrais manger, boire ou parler avec une gueule aussi endiablée ? Autant me gaver et me soûler de diables à même le plat ou la cruche, puisque je me rendrais assurément complice de tous les diables qui habitent dans les Juifs et conspuent le sang précieux du Christ. Dieu m’en garde !}
  \chunks{\textelp{}}
         {\textelp{}}
% WA 53, p. 528-529
  \chunks{Insonderheit, wo ihr Prediger seid, da Juden sind, da halt an mit Fleiß bei euren Herren und Regenten, daß sie ihr Amt bedenken, wie sie Gott schuldig sind, und die Juden zur Arbeit zwingen, den Wucher verbieten und steuren ihrem Lästern und Fluchen. Denn so sie unter uns Christen die Diebe, Räuber, Mörder, Lästerer und ander Laster strafen, warum sollen die Juden Teufels frei sein, solchs bei und wider uns zu üben? Leiden wir doch mehr von ihnen, weder die Walhen von den Spaniolen. Dieselbigen nehmen dem Hauswirt Küchen, Keller, Kasten, Beutel ein, fluchen ihnen dazu und dräuen ihnen den Tod. Also tun uns die Juden, unser Gäste, auch; wir sind ihre Hauswirte.}
         {En particulier, là où vous êtes prédicateurs et où il y a des Juifs, insistez avec zèle auprès de vos seigneurs et gouvernants pour qu’ils considèrent leur charge et ce qu’ils doivent à Dieu, qu’ils contraignent les Juifs au travail, interdisent l’usure et mettent un frein à leurs blasphèmes et à leurs malédictions. Car s’ils punissent parmi nous, chrétiens, les voleurs, les brigands, les meurtriers, les blasphémateurs et les autres vices, pourquoi diable les Juifs seraient-ils libres de les pratiquer chez nous et contre nous ? Nous souffrons pourtant d’eux plus que les Italiens des Espagnols\footnote{Allusion probable aux soldats espagnols logés chez l’habitant pendant les guerres d’Italie.}. Ceux-ci mettent à sac la cuisine, la cave, le coffre et la bourse de leur logeur, le maudissent par-dessus le marché et le menacent de mort. Ainsi nous font aussi les Juifs, nos pensionnaires : nous sommes leurs logeurs.}
% WA 53, p. 529
  \chunks{So rauben sie und saugen uns aus, liegen uns auf dem Halse, die faulen Schelmen und müßigen Wänste, saufen, fressen, haben gute Tage in unserm Hause, fluchen zu Lohn unserm Herrn Christo, Kirchen, Fürsten und uns allen, dräuen und wünschen uns ohn Unterlaß den Tod und alles Unglück. Denke doch, wo kommen wir armen Christen dazu, daß wir solch faul, müßig Volk, solch unnütz, böse, schädlich Volk, solche lästerliche Feinde Gottes umsonst sollen nähren und reich machen, dafür nichts kriegen denn ihr Fluchen, Lästern und alles Unglück, das sie uns tun und wünschen können? Sind wir doch wohl so blind und starrende Klötze in diesem Stücke als die Juden in ihrem Unglauben, daß wir solch große Tyrannei von den heillosen Schelmen leiden, solchs nicht sehen noch fühlen, wie sie unser Junker, ja unsere wütige Tyrannen sind, wir aber ihre Gefangene und Untertan; klagen noch, sie seien unsere Gefangene, spotten unser dazu, als müßten wirs von ihnen leiden.}
         {Eux, au contraire, nous pillent et nous sucent, nous sont à charge, ces coquins fainéants et ces panses oisives : ils bâfrent, s’enivrent, se donnent du bon temps dans notre maison, et pour salaire maudissent notre Seigneur le Christ, les Églises, les princes et nous tous, nous menacent et nous souhaitent sans relâche la mort et tous les malheurs. Pense donc : comment en venons-nous, pauvres chrétiens, à devoir nourrir et enrichir pour rien un peuple si paresseux, si oisif, un peuple si inutile, si méchant, si nuisible, de tels ennemis blasphémateurs de Dieu, sans recevoir en échange autre chose que leurs malédictions, leurs blasphèmes et tous les malheurs qu’ils peuvent nous faire et nous souhaiter ? Nous sommes bien, en cette affaire, des souches aussi aveugles et insensibles que les Juifs dans leur incrédulité, pour souffrir une si grande tyrannie de ces scélérats sans la voir ni la sentir : ils sont nos hobereaux, et même nos tyrans furieux, et nous leurs captifs et leurs sujets ; et ils se plaignent encore d’être nos captifs, et se moquent de nous par-dessus le marché, comme s’il nous fallait le souffrir d’eux.}
% WA 53, p. 529
  \chunks{Wollen aber die Herren sie nicht zwingen noch solchem ihrem teuflischen Mutwillen steuren, daß man sie zum Lande austreibe, wie gesagt, und lassen ihnen sagen, daß sie hinziehen in ihr Land und Güter gen Jerusalem und daselbs lügen, fluchen, lästern, speien, morden, stehlen, rauben, wuchern, spotten und alle solche lästerliche Greuel treiben, wie sie bei uns tun, und lassen uns unser Herrschaft, Land, Leib und Gut, viel mehr unsern Herrn Messiam, Glauben und Kirchen unbeschweret und unbeschmeisset mit solchen ihren teuflischen Tyranneien und Bosheiten. Ob sie Freiheit fürwenden könnten, die sollen ihnen nichts helfen; denn es kann niemand Freiheit geben, solche Greuel zu üben, und sind alle Freiheit dadurch verwahrloset und verloren.}
         {Mais si les seigneurs ne veulent pas les contraindre ni mettre un frein à leur malice diabolique, qu’on les chasse du pays, comme on l’a dit, et qu’on leur fasse dire de s’en aller dans leur pays et dans leurs biens, à Jérusalem, et d’y mentir, maudire, blasphémer, conspuer, assassiner, voler, piller, pratiquer l’usure, se moquer et commettre toutes ces abominations blasphématoires, comme ils le font chez nous ; et qu’ils nous laissent notre seigneurie, notre pays, nos corps et nos biens, et bien plus encore notre Seigneur le Messie, notre foi et nos Églises, sans les accabler ni les souiller de leurs tyrannies et méchancetés diaboliques. S’ils pouvaient alléguer des privilèges, ceux-ci ne devront leur servir de rien ; car personne ne peut accorder le privilège de commettre de telles abominations, et par là tous les privilèges sont déchus et perdus.}
% WA 53, p. 529
  \chunks{Wenn ihr Pfarrherren und Prediger (neben mir) solch treulich Warnen habt ausgericht und will weder Herr noch Untertan etwas dazu tun, so laßt uns (wie Christus spricht) den Staub von den Schuhen schütteln und sagen: Wir sind unschuldig an eurem Blut. Denn ich sehe wohl und habs oft erfahren, wie gar barmherzig die verkehrete Welt ist, da sie billich sollt scharf sein, und widerum scharf ist, da sie sollt barmherzig sein, wie der König Ahab, 3. Reg. 20. So regiert der Fürst dieser Welt. Also werden sie vielleicht itzt auch barmherzig sein wollen über die Juden, die blutdürstigen Feinde unsers christlichen und menschlichen Namens, damit den Himmel zu verdienen.}
         {Si vous, pasteurs et prédicateurs, avez avec moi loyalement donné un tel avertissement, et que ni seigneur ni sujet ne veuille rien faire, alors secouons (comme dit le Christ) la poussière de nos souliers, et disons : Nous sommes innocents de votre sang. Car je vois bien, et je l’ai souvent éprouvé, combien le monde à l’envers est miséricordieux là où il devrait à bon droit être âpre, et à son tour âpre là où il devrait être miséricordieux, comme le roi Achab, 1~Rois~20\footnote{En 1~Rois~20, Achab épargne Ben-Hadad, roi d’Aram, que Dieu avait voué à l’anathème, et un prophète le lui reproche (v.~42). Le \enquote{prince de ce monde} est le diable (Jean~12,~31).}. Ainsi gouverne le prince de ce monde. De même voudront-ils peut-être maintenant être miséricordieux envers les Juifs, les ennemis assoiffés de sang de notre nom chrétien et humain, pour mériter par là le ciel.}
% WA 53, p. 530
  \chunks{Aber daß die Juden mit allen solchen teuflischen obererzählten Greueln uns arme Christen fangen, plagen, martern und alle Herzeleid anlegen, das soll man leiden, und ist christlich wohl getan, sonderlich so Geld da ist, das sie uns gestohlen und geraubt haben.}
         {Mais que les Juifs, avec toutes les abominations diaboliques énumérées plus haut, nous capturent, nous tourmentent, nous martyrisent, nous, pauvres chrétiens, et nous infligent toute sorte de chagrin, cela, il faut le souffrir, et c’est agir en bon chrétien, surtout s’il y a là de l’argent qu’ils nous ont volé et pillé.}
  \chunks{\textelp{}}
         {\textelp{}}
% WA 53, p. 530
  \chunks{Ich hab viel Historien gelesen und gehört von den Juden, so mit diesem Urteil Christi stimmen, nämlich wie sie die Brunnen vergiftet, heimlich gemordet, Kinder gestohlen, wie droben gemeldet. Item, daß ein Jude dem andern über Feld einen Topf voll Bluts, auch durch einen Christen, zugeschickt; item, ein Faß Wein, da das ausgetrunken, ein toter Jude im Fasse gefunden; und dergleichen viel. Und das Kinderstehlen hat sie oft (wie droben gesagt) verbrennet und verjagt. Ich weiß wohl, daß sie solches und alles leugnen. Es stimmet aber alles mit dem Urteil Christi, daß sie giftige, bittere, rachgierige, hämische Schlangen, Meuchelmörder und Teufelskinder sind, die heimlich stechen und Schaden tun, weil sie es öffentlich nicht vermögen. Darum ich gerne wollte, sie wären, da keine Christen sind. Der Türke und ander Heiden leiden solchs nicht von ihnen, das wir Christen von den giftigen Schlangen und jungen Teufeln leiden; sie tuns auch niemand denn uns Christen. Das ists, das ich droben gesagt habe, daß ein Christ nächst dem Teufel keinen giftigern, bittern Feind habe denn einen Juden, so wir doch niemand so viel Guts tun noch so viel von jemand leiden als eben von solchen bösen Teufelskindern und Schlangengezüchte.}
         {J’ai lu et entendu raconter sur les Juifs beaucoup d’histoires qui s’accordent avec ce jugement du Christ\footnote{Matthieu~12,~34 et 3,~7 (\enquote{engeance de vipères}), cités au paragraphe précédent.} : comment ils ont empoisonné les puits, assassiné en secret, volé des enfants, comme on l’a rapporté plus haut. De même, qu’un Juif a envoyé à un autre, à travers la campagne, un pot plein de sang, et cela par l’intermédiaire d’un chrétien ; de même, un tonneau de vin au fond duquel, une fois bu, on a trouvé un Juif mort ; et bien d’autres choses semblables. Et le vol d’enfants leur a souvent valu (comme on l’a dit plus haut) d’être brûlés et chassés. Je sais bien qu’ils nient cela et tout le reste. Mais tout s’accorde avec le jugement du Christ : ce sont des serpents venimeux, amers, vindicatifs, perfides, des assassins sournois et des enfants du diable, qui piquent et nuisent en secret, tant qu’ils ne le peuvent pas ouvertement. C’est pourquoi je voudrais bien qu’ils soient là où il n’y a pas de chrétiens. Le Turc et les autres païens ne souffrent pas d’eux ce que nous, chrétiens, souffrons de ces serpents venimeux et de ces jeunes diables ; aussi bien ne le font-ils à personne d’autre qu’à nous, chrétiens. C’est ce que j’ai dit plus haut : un chrétien n’a pas, après le diable, d’ennemi plus venimeux et plus amer qu’un Juif ; et pourtant nous ne faisons à personne autant de bien, et ne souffrons de personne autant que de ces méchants enfants du diable et de cette engeance de serpents.}
  \chunks{\textelp{}}
         {\textelp{}}
% WA 53, p. 538
  \chunks{Ja, wie wollen wir tun? Wenn wir gleich den Juden ihre Synagoga verbrennen, Gott loben, beten, lehren, Gottes Namen nennen öffentlich verbieten etc., gleichwohl werden sie es doch heimlich nicht lassen. Und weil wir wissen, daß sie es heimlich tun, so ists eben so viel, als täten sie es öffentlich. Denn was man weiß, daß heimlich geschicht und geduldet wird, das heißt doch nicht heimlich, und gleichwohl unser Gewissen damit vor Gott beschweret ist. Wohlan, da mögen wir uns fürsehen.}
         {Oui, mais que ferons-nous ? Quand bien même nous brûlerions les synagogues des Juifs et leur interdirions de louer Dieu, de prier, d’enseigner, de prononcer publiquement le nom de Dieu, etc., ils ne laisseront pas pour autant de le faire en secret. Et puisque nous savons qu’ils le font en secret, c’est exactement comme s’ils le faisaient publiquement. Car ce qu’on sait se faire en secret, et qu’on tolère, ne s’appelle plus secret, et notre conscience n’en est pas moins chargée devant Dieu. Eh bien, prenons-y garde.}
% WA 53, p. 538
  \chunks{Meins Dünkens wills doch da hinaus: Sollen wir der Juden Lästerung rein bleiben und nicht teilhaftig werden, so müssen wir gescheiden sein, und sie aus unserm Lande vertrieben werden. Sie mögen gedenken in ihr Vaterland; so dürfen sie nicht mehr vor Gott über uns schreien und lügen, daß wir sie gefangen halten, wir auch nicht klagen, daß sie uns mit ihrem Lästern und Wuchern beschweren. Dies ist der nähest und beste Rat, der beide Part in solchem Fall sichert.}
         {À mon sens, cela revient pourtant à ceci : si nous voulons rester purs des blasphèmes des Juifs et n’y avoir aucune part, il faut que nous soyons séparés, et qu’ils soient chassés de notre pays. Qu’ils songent à retourner dans leur patrie ; alors ils n’auront plus besoin de crier contre nous devant Dieu et de mentir en disant que nous les tenons captifs, et nous n’aurons plus à nous plaindre qu’ils nous accablent de leurs blasphèmes et de leur usure. Voilà le conseil le plus simple et le meilleur, qui met en pareil cas les deux parties en sûreté.}
% WA 53, p. 538
  \chunks{Aber hie werden sie, als die das Land ungern räumen, getrost alles und alles leugnen, dazu auch Gelds genug der Herrschaft bieten, ob sie bleiben möchten.}
         {Mais ici, en gens qui quittent le pays à contrecœur, ils nieront hardiment tout, absolument tout, et offriront en outre aux seigneurs assez d’argent, dans l’espoir de pouvoir rester.}
% WA 53, p. 538
  \chunks{Wehe aber denen, so solch Geld nehmen, und verflucht sei solch Geld, das sie doch auch sonst verfluchter Weise durch Wucher uns gestohlen haben. Denn sie leugnen ja so sehr, als sie lügen, und wo sie uns Christen heimlich fluchen, giften oder Schaden tun können, des machen sie ihnen kein Gewissen. Werden sie darüber begriffen oder sonst bezichtiget, so türen sie es getrost leugnen, auch bis in den Tod, weil sie uns nicht wert achten, denen sie sollten die Wahrheit bekennen, nachdem es die heiligen Gotteskinder gewiß dafür halten, daß, was sie uns Böses fluchen und tun können, einen großen Gottesdienst daran tun. Ja, wenn sie uns das könnten tun, das wir ihnen tun können, würde unser keiner eine Stunde leben müssen. Weil sie es aber öffentlich nicht vermögen zu tun, bleiben sie gleichwohl im Herzen unser tägliche Mörder und blutdürstige Feinde.}
         {Mais malheur à ceux qui prennent un tel argent, et maudit soit cet argent, qu’ils nous ont d’ailleurs volé de façon maudite par l’usure ! Car ils sont aussi prompts à nier qu’à mentir ; et là où ils peuvent secrètement nous maudire, nous empoisonner ou nous nuire, à nous chrétiens, ils ne s’en font pas scrupule. S’ils sont pris sur le fait ou accusés de quelque autre manière, ils osent le nier hardiment, jusqu’à la mort même\footnote{Sous la torture, selon la glose de WA.}, parce qu’ils ne nous jugent pas dignes qu’on nous avoue la vérité, puisque les saints enfants de Dieu tiennent assurément pour un grand service de Dieu tout le mal qu’ils peuvent nous souhaiter et nous faire. Oui, s’ils pouvaient nous faire ce que nous pouvons leur faire, pas un de nous ne resterait en vie une heure. Mais puisqu’ils ne peuvent le faire ouvertement, ils n’en restent pas moins, dans leur cœur, nos meurtriers de tous les jours et nos ennemis assoiffés de sang.}
% WA 53, p. 538
  \chunks{Solchs beweisen ihr Beten und Fluchen und so viel Historien, da sie Kinder gemartert und allerlei Laster geübt, darüber sie oft verbrannt und verjagt sind.}
         {C’est ce que prouvent leurs prières et leurs malédictions, et tant d’histoires où ils ont martyrisé des enfants et commis toutes sortes de crimes, ce qui les a souvent fait brûler et chasser.}
% WA 53, p. 538
  \chunks{Darum ich wohl glaube, daß sie viel Ärgers heimlich reden und tun, weder die Historien und andere von ihnen schreiben; sie aber sich aufs Leugnen und ihr Geld verlassen.}
         {C’est pourquoi je crois bien qu’ils disent et font en secret bien pire que ce qu’en écrivent les histoires et d’autres auteurs ; mais ils se fient à leurs dénégations et à leur argent.}
% WA 53, p. 538-539
  \chunks{Aber wenn sie gleich alles leugnen könnten, so können sie das nicht leugnen, daß sie uns Christen öffentlich fluchen, nicht um unser böses Leben willen, sondern daß wir den Jesum für Messia halten, und daß sie bei uns gefangen sein müssen, so sie wohl wissen, daß sie daran lügen, und viel mehr sie uns in unserm Lande durch ihren Wucher gefangen halten, jedermann aber ihr gern los wäre.}
         {Mais quand bien même ils pourraient tout nier, ils ne peuvent nier qu’ils nous maudissent publiquement, nous chrétiens, non à cause de notre mauvaise vie, mais parce que nous tenons ce Jésus pour le Messie ; ni qu’ils se disent captifs chez nous, alors qu’ils savent bien qu’ils mentent en cela, et que ce sont bien plutôt eux qui nous tiennent captifs dans notre pays par leur usure, tandis que chacun serait heureux d’en être débarrassé.}
% WA 53, p. 539
  \chunks{Weil sie aber uns verfluchen, so verfluchen sie unsern Herrn auch; verfluchen sie unsern Herrn, so verfluchen sie auch Gott den Vater, Schöpfer Himmels und der Erden, daß also ihr Leugnen sie nichts helfen kann. Sie sind allein mit dem Fluchen überwunden, daß man alles wohl glauben muß, was man Böses von ihnen schreibet; sie tuns gewißlich mehr und ärger, weder wir wissen oder erfahren. Denn Christus leugt und treugt nicht, der sie Schlangen und Teufelskinder urteilt, das ist: seine und aller der Seinen Mörder und Feinde, wo sie können.}
         {Or, puisqu’ils nous maudissent, ils maudissent aussi notre Seigneur ; et s’ils maudissent notre Seigneur, ils maudissent aussi Dieu le Père, créateur du ciel et de la terre ; si bien que leurs dénégations ne peuvent leur servir de rien. Par leurs seules malédictions, ils sont convaincus, de sorte qu’il faut bien croire tout le mal qu’on écrit d’eux ; ils en font assurément davantage et de pires que nous n’en savons ou n’en apprenons. Car le Christ ne ment ni ne trompe, lui qui les juge serpents et enfants du diable, c’est-à-dire ses meurtriers et ses ennemis, et ceux de tous les siens, là où ils le peuvent.}
  \chunks{\textelp{}}
         {\textelp{}}
% WA 53, p. 541-542
  \chunks{Unsern Oberherren, so Juden unter sich haben, wünsche ich und bitte, daß sie eine scharfe Barmherzigkeit wollten gegen diese elende Leute üben, wie droben gesagt, obs doch etwas (wiewohl es mißlich ist) helfen wollte, wie die treuen Ärzte tun, wenn das heilige Feuer in die Bein kommen ist: fahren sie mit Unbarmherzigkeit und schneiden, sägen, brennen Fleisch, Adern, Bein und Mark abe. Also tu man hie auch: verbrenne ihr Synagogen, verbiete alles, was ich droben erzählet habe, zwinge sie zur Arbeit und gehe mit ihnen um nach aller Unbarmherzigkeit, wie Mose tät in der Wüsten und schlug dreitausend tot, daß nicht der ganze Haufe verderben mußte. Sie wissen wahrlich nicht, was sie tun, wollens dazu, wie die besessen Leute, nicht wissen, hören noch lernen. Darum kann man hie keine Barmherzigkeit üben, sie in ihrem Wesen zu stärken. Will das nicht helfen, so müssen wir sie wie die tollen Hunde ausjagen, damit wir nicht, ihrer greulichen Lästerung und aller Laster teilhaftig, mit ihnen Gottes Zorn verdienen und verdammt werden. Ich habe das Meine getan; ein jeglicher sehe, wie er das Seine tu. Ich bin entschuldigt.}
         {À nos souverains qui ont des Juifs sous leur autorité, je souhaite et je demande qu’ils veuillent exercer envers ces misérables une âpre miséricorde, comme on l’a dit plus haut, pour voir si cela pourrait servir à quelque chose (bien que ce soit douteux) ; comme font les médecins fidèles quand le feu Saint-Antoine s’est mis dans les os : ils procèdent sans miséricorde, et coupent, scient, brûlent la chair, les veines, les os et la moelle. Qu’on fasse de même ici : qu’on brûle leurs synagogues, qu’on interdise tout ce que j’ai énuméré plus haut, qu’on les contraigne au travail, et qu’on agisse avec eux sans aucune miséricorde, comme Moïse fit au désert, lorsqu’il en tua trois mille, afin que la multitude entière ne pérît pas. Ils ne savent vraiment pas ce qu’ils font, et de plus, comme des possédés, ils ne veulent ni le savoir, ni l’entendre, ni l’apprendre. C’est pourquoi on ne peut pas exercer ici une miséricorde qui les affermirait dans leurs agissements. Si cela ne sert à rien, il nous faudra les chasser comme des chiens enragés, afin que nous ne méritions pas avec eux la colère de Dieu et ne soyons pas damnés, en prenant part à leurs abominables blasphèmes et à tous leurs vices. J’ai fait ce qui me revient ; que chacun voie comment il fera ce qui lui revient. Je suis déchargé.}

  \chunks{\begin{flushright}
          \large\textsc{Luther}, \textit{Von den Juden und ihren Lügen} (1543).
          \end{flushright}}
         {\begin{flushright}
          \large\textsc{Luther}, \textit{Des Juifs et de leurs mensonges}.\par
          \normalsize Traduction nouvelle.
          \end{flushright}}
\end{paracol}

\input{glossaire.inc}
\end{document}
